% !TeX root = main.tex
% !TeX spellcheck = pt_BR

\chapter{Inferência}


\section{Distribuições amostrais}

\begin{theorem}
	Sejam $X_1, X_2, \ldots, X_n$ uma A.A. e seja $g$ uma função real tal que $\E[g(X)]$ e $\Var[g(X)]$ existam. Nesse caso,
	\[
	\E\left[\sum_{i=1}^{n} g(X_i)\right] = n\E[g(X)],
	\]
	e
	\[
	\Var\left[\sum_{i=1}^{n} g(X_i)\right] = n\Var[g(X)].
	\]
\end{theorem}

\begin{theorem}
	Sejam $X_1, X_2, \ldots, X_n$ uma A.A. de uma população com média $\mu$ e variância $\sigma^2 < \infty$. Então,
	\begin{outline}[enumerate]
		\1 $\E(\bar{X}) = \mu$;
		
		\1 $\Var(\bar{X}) = \sigma^2/n$;
		
		\1 $\E(S^2) = \sigma^2$.
	\end{outline}
\end{theorem}

\begin{theorem}[Distribuição amostral da média]
	Sejam $X_1, X_2, \ldots, X_n$ uma A.A. de uma população com função geradora de momentos $M_X(t)$. Então, a FGM da média amostral será
	\[
	M_{\bar{X}}(t) = [M_X(t/n)]^n.
	\]
\end{theorem}

\begin{theorem}[A.A. da Normal]
	Sejam $X_1, X_2, \ldots, X_n$ uma A.A. de uma população com distribuição $\normdist{\mu, \sigma^2}$. Então,
	\begin{outline}[enumerate]
		\1 $\bar{X}$ e $S^2$ são independentes;
		
		\1 $\bar{X} \sim \normdist{\mu, \frac{\sigma^2}{n}}$;
		
		\1 $\frac{(n-1)S^2}{\sigma^2} \sim \chi_{n-1}^2$.
	\end{outline}
\end{theorem}

\begin{theorem}[A.A. da t-Student]
	Sejam $X_1, X_2, \ldots, X_n$ uma A.A. oriunda de uma população $\normdist{\mu, \sigma^2}$. Então
	\[
	\frac{\bar{X} - \mu}{S/\sqrt{n}} \sim \tdist{n-1}.
	\]
	\begin{note}
		se $Z \sim \normdist{0,1}$ e $V \sim \chi_k^2$ independentes, então
		\[
		\frac{Z}{\sqrt{V/k}} \sim \tdist{k}.
		\]
	\end{note}
\end{theorem}

\begin{theorem}[Razão de A.A. da Normal]
	Sejam $X_1, X_2, \ldots, X_n$ e $Y_1, Y_2, \ldots, Y_m$ duas A.A.'s oriundas de populações $\normdist{\mu_X, \sigma_X^2}$ e $\normdist{\mu_Y, \sigma_Y^2}$, respectivamente. Então
	\[
	\frac{S_X^2/\sigma_X^2}{S_Y^2/\sigma_Y^2} \sim F_{n-1, m-1}.
	\]
	\begin{note}
		O resultado acima decorre do fato de que, se $V \sim \chi_r^2$ e $W \sim \chi_s^2$ independentes, então
		\[
		\frac{V/r}{W/s} \sim F_{r,s}.
		\]
	\end{note}
\end{theorem}

\subsection{Estatísticas de ordem}


\begin{theorem}[Máximo e mínimo]
	\leavevmode
	\begin{outline}[enumerate]
		\1 \( F_{X_{(n)}}(x) = \left[ F_X(x) \right]^n \) (contínua e discreta)
		\1 \( f_{X_{(n)}}(x) = n\left[ F_X(x) \right]^{n-1} f_X(x) \) (contínua)
		\1 \( F_{X_{(1)}}(x) = 1-\left[1 - F_X(x) \right]^n \) (contínua e discreta)
		\1 \( f_{X_{(1)}}(x) = n\left[1- F_X(x) \right]^{n-1} f_X(x) \) (contínua)
	\end{outline}
	
\end{theorem}

\begin{theorem}[Contínua e discreta]
	Sejam $X_1, X_2, \ldots, X_n$ uma A.A. de uma população com f.d. $F_X$. Nesse caso, a f.d. da j-ésima estatística de ordem, $X_{(j)}$, será
	\[
	F_{X_{(j)}}(x) = \sum_{k=j}^{n} \binom{n}{k} [F_X(x)]^k [1 - F_X(x)]^{n-k}.
	\]
\end{theorem}

\begin{theorem}[Discreta]
	Sejam $X_1, X_2, \ldots, X_n$ uma A.A. de uma população com f.d. $F_X$. Suponha que a população seja discreta com f.p. $f_X$ e suporte $\{x_i \mid x_1 < x_2 < \ldots < x_n\}$. Além disso, denote $P_0 = 0$ e $P_i = F_X(x_i) = f_X(x_1) + f_X(x_2) + \ldots + f_X(x_i)$, $i = 1, 2, \ldots, n$. Nesse caso, a função de probabilidade da j-ésima estatística de ordem será
	\[
	f_{X_{(j)}}(x) = \P(X_{(j)} = x_i) = \sum_{k=j}^{n} \binom{n}{k} \left[ P_i^k (1 - P_i^{n-k}) - P_{i-1}^k (1 - P_{i-1}^{n-k}) \right].
	\]
\end{theorem}

\begin{theorem}[Contínua]
	Sejam $X_1, X_2, \ldots, X_n$ uma A.A. de uma população com f.d. $F_X$. Se a população for contínua com f.d.p. $f_X$, então a f.d.p. da j-ésima estatística de ordem $X_{(j)}$ será
	\[
	f_{X_{(j)}}(x) = \frac{n!}{(j-1)!(n-j)!} f_X(x) [F_X(x)]^{j-1} [1 - F_X(x)]^{n-j}.
	\]
\end{theorem}

\begin{theorem}[Contínua]
	Sejam $X_1, X_2, \ldots, X_n$ uma A.A. de uma população com f.d. $F_X$ e f.d.p. $f_X$. Então a f.d.p. conjunta de $(X_{(i)}, X_{(j)})$, $1 \leq i < j \leq n$, será
	\begin{align*}
		f_{X_{(i)}, X_{(j)}}(u, v) &= \frac{n!}{(i-1)!(j-i-1)!(n-j)!} f_X(u) f_X(v) [F_X(u)]^{i-1} \\
		&\quad \times [F_X(v) - F_X(u)]^{j-i-1} [1 - F_X(v)]^{n-j} \ind(u < v).
	\end{align*}
	\begin{note}
		Argumentos análogos podem ser utilizados para derivar distribuições de v.a.'s com três ou mais estatísticas de ordem. Um dos v.a.'s mais utilizados corresponde a $(X_{(1)}, X_{(2)}, \ldots, X_{(n)})$, cuja f.d.p. corresponde a
		\[
		f_{X_{(1)}, \ldots, X_{(n)}}(x_1, \ldots, x_n) = n! f_X(x_1) \cdots f_X(x_n) \ind(-\infty < x_1 < \ldots < x_n < \infty).
		\]
	\end{note}
\end{theorem}

\section{Redução de dados}

\begin{definition}[Princípio da suficiência]
	Se $T$ for uma estatística suficiente para $\theta$, e se duas amostras observadas, $\v{x}$ e $\v{y}$, forem tais que $T(\v{x}) = T(\v{y})$, então as inferências a respeito de $\theta$ deverão ser as mesmas, independentemente de se usar $\v{x}$ ou $\v{y}$.
	
	\begin{note}
		Como consequência, qualquer inferência sobre $\theta$ deverá depender da amostra $\v{X}$ apenas através de $T(\v{X})$.
	\end{note}
\end{definition}

\begin{definition}[Estatística suficiente]
	Uma estatística $T(\v{X})$ é chamada de estatística suficiente para $\theta$ se a distribuição condicional da amostra $\v{X}$, dado o valor de $T(\v{X})$, não depender de $\theta$.
\end{definition}

\begin{theorem}[Razão para suficiência]
	Seja $p(\v{x} | \theta)$ a f.p. (ou f.d.p.) conjunta de uma amostra $\v{X}$, e $q(t | \theta)$ a f.p. (ou f.d.p.) de uma estatística $T(\v{X})$. Então $T(\v{X})$ será uma estatística suficiente para $\theta$ se, para todo $\v{x} \in \mathcal{X}$, a razão
	\[
	\frac{p(\v{x} | \theta)}{q(T(\v{x}) | \theta)}
	\]
	for constante como função de $\theta$.
\end{theorem}

\begin{theorem}[Critério da fatoração]
	Seja $f(\v{x} | \theta)$ a f.p. (ou f.d.p.) conjunta de uma amostra $\v{X}$. Uma estatística $T(\v{X})$ é uma estatística suficiente para $\theta$ se, e somente se, existirem funções $g(t | \theta)$ e $h(\v{x})$ de modo que, para todo $\v{x} \in \mathcal{X}$ e todo $\theta \in \Theta$,
	\[
	f(\v{x} | \theta) = g(T(\v{x}) | \theta)h(\v{x}).
	\]
\end{theorem}

\begin{theorem}[Suficiência na família exponencial]
	Sejam $X_1, X_2, \ldots, X_n$ uma A.A. oriunda de uma população cuja f.d.p. (ou f.p.) $f(x | \theta)$ pertence à família exponencial
	\[
	f(x | \v{\theta}) = \exp\left(\sum_{k=1}^{s} \eta_k(\v{\theta})t_k(x)\right) a(\v{\theta})h(x),
	\]
	em que $\v{\theta} = (\theta_1, \theta_2, \ldots, \theta_d)$, $d \leq s$. Então,
	\[
	T(\v{X}) = \left(\sum_{i=1}^{n} t_1(X_i), \sum_{i=1}^{n} t_2(X_i), \ldots, \sum_{i=1}^{n} t_s(X_i)\right)
	\]
	é uma estatística suficiente para $\v{\theta}$.
\end{theorem}

\begin{definition}[Estatística suficiente mínima]
	Uma estatística $T(\v{X})$ é chamada estatística suficiente mínima se, para qualquer outra estatística suficiente $T'(\v{X})$, $T(\v{X})$ puder ser escrita como função de $T'(\v{X})$.
\end{definition}

\begin{theorem}[Critério para suficiência mínima]
	Seja $f(\v{x} | \theta)$ a f.p. (ou f.d.p.) conjunta de uma amostra $\v{X}$. Suponhamos que, para quaisquer duas possíveis amostras observadas, $\v{x}$ e $\v{y}$, $f(\v{x} | \theta)/f(\v{y} | \theta)$ seja constante como função de $\theta$ se, e somente se, $T(\v{x}) = T(\v{y})$. Então $T(\v{X})$ é uma estatística suficiente mínima.
\end{theorem}

\begin{definition}[Estatística ancilar]
	Uma estatística $S(\v{X})$ cuja distribuição não depende de $\theta$ é chamada de estatística ancilar.
	
	\begin{note}
		Uma estatística $S(\v{X})$ é dita ser ancilar de primeira ordem se $\E_\theta[S(\v{X})]$ não depende de $\theta$.
	\end{note}
\end{definition}

\begin{definition}[Estatística completa]
	Seja $f(t | \theta)$ uma família de f.d.p.'s (ou f.p.'s) para uma estatística $T(\v{X})$. A família de distribuições de probabilidade é chamada completa se, para qualquer função $g$,
	\[
	\E_\theta[g(T)] = 0 \implies \P_\theta[g(T) = 0] = 1,
	\]
	para todo $\theta \in \Theta$. De forma equivalente, $T(\v{X})$ é chamada de estatística completa.
\end{definition}

\begin{theorem}[Basu]
	Se $T(\v{X})$ for uma estatística suficiente mínima e completa, então $T(\v{X})$ será independente de toda estatística ancilar.
\end{theorem}

\begin{theorem}[Estatísticas completas na família exponencial]
	Sejam $X_1, X_2, \ldots, X_n$ uma A.A. oriunda de uma população com f.d.p. (ou f.p.) pertencente à família exponencial
	\[
		f(x | \v{\theta}) = \exp\left(\sum_{k=1}^{s} \eta_k(\v{\theta})t_k(x)\right) a(\v{\theta})h(x),
	\]
	em que $\v{\theta} = (\theta_1, \theta_2, \ldots, \theta_s)$. Então, a estatística
	\[
		T(\v{X}) = \left(\sum_{i=1}^{n} t_1(X_i), \sum_{i=1}^{n} t_2(X_i), \ldots, \sum_{i=1}^{n} t_s(X_i)\right)
	\]
	é completa, desde que $\Theta$ contenha um conjunto aberto em $\Re^s$.
\end{theorem}

\begin{theorem}
	Se uma estatística suficiente mínima existe, então qualquer estatística completa também será uma estatística suficiente mínima.
\end{theorem}

\section{Estimação Pontual I}

\begin{definition}[Estimador pontual]
	Um estimador pontual é qualquer função de uma amostra.
	
	\begin{note}
		Qualquer estatística é um estimador.
	\end{note}
\end{definition}

\begin{definition}[Estimativa]
	Uma estimativa é um valor (valores) observado(s) de um estimador, obtido a partir de uma amostra observada.
\end{definition}

\begin{definition}[Função de verossimilhança]
	Sejam $X_1, X_2, \ldots, X_n$ uma amostra oriunda de uma população com f.d.p. (ou f.p.) $f(x | \theta)$. A função de verossimilhança será
	\[
	L(\theta | \v{x}) = f(\v{x} | \theta).
	\]
	Qualquer estimador $\hat{\theta} \equiv \hat{\theta}_n(\v{X}) \in \Theta$ tal que
	\[
	L(\hat{\theta} | \v{X}) = \sup_{\theta \in \Theta} L(\theta | \v{X})
	\]
	é chamado estimador de máxima verossimilhança (EMV) de $\theta$.
\end{definition}

\begin{theorem}[Pontos críticos]
	Seja $g(x)$ uma função contínua e duplamente diferenciável em algum intervalo aberto $I$ que contém $a$. Se $g'(a) = 0$, então $g(x)$ tem um máximo (mínimo) relativo em $x = a$, desde que $g''(a) < 0$ ($g''(a) > 0$). Ademais, o conjunto $\{a \in I \mid g'(a) = 0\}$ é chamado de conjunto dos pontos críticos de $g(x)$. Os únicos máximos (mínimos) possíveis de $g(x)$ para $c \leq x \leq d$ estão nos pontos críticos entre $c$ e $d$, bem como em $c$ e em $d$.
\end{theorem}

\begin{theorem}[Propriedade da invariância]
	Seja $\hat{\theta}$ o EMV de $\theta$. Então, para qualquer estimando $\tau(\theta)$, seu EMV será $\tau(\hat{\theta})$.
\end{theorem}

\begin{theorem}[Relação EMV e Suficiência]
	Seja $L(\theta | \v{x})$ a função de verossimilhança da amostra $X_1, X_2, \ldots, X_n$, e seja $T = T(\v{X})$ uma estatística suficiente para $\theta$. O EMV de $\theta$, se existir, será função de $T$.
\end{theorem}

\begin{theorem}[EMV na família exponencial]
	Sejam $X_1, X_2, \ldots, X_n$ uma A.A. pertencente à família exponencial canônica, cuja f.d.p. (ou f.p.) conjunta pode ser escrita como
	\[
	f(\v{x} | \eta) = \exp\left(\eta \sum_{i=1}^{n} T(x_i) - nA(\eta)\right) h^*(\v{x}).
	\]
	Temos que o EMV de $\eta$ será solução da equação
	\[
	\frac{1}{n} \sum_{i=1}^{n} T(X_i) = \E[T(X_1)],
	\]
	desde que essa solução pertença ao espaço paramétrico de $\eta$.
	
	\begin{note}
		Devido à propriedade da invariância, o resultado permanece válido para o caso geral da família exponencial, em que cada $X_i$ tem f.d.p. (ou f.p.)
		\[
		f(x | \theta) = \exp\{\eta(\theta)T(x) - B(\theta)\} h(x).
		\]
	\end{note}
	
	\begin{note}
		O Teorema 4 pode ser generalizado para o caso multiparamétrico:
		\[
		f(\v{x} | \v{\theta}) = \exp\left(\sum_{k=1}^{s} \eta_k(\v{\theta})T_k(\v{x}) - B(\v{\theta})\right) h(\v{x}).
		\]
	\end{note}
\end{theorem}

\section{Estimação Pontual II}

\begin{definition}[Erro quadrático médio]
	O erro quadrático médio (EQM) do estimador $W$ de um parâmetro $\theta$ corresponde a
	\[
	\E_\theta(W - \theta)^2,
	\]
	que é função de $\theta$.
\end{definition}

\begin{definition}[Viés]
	O viés ($B$ -- do inglês, bias) do estimador $W$ de um parâmetro $\theta$ corresponde a
	\[
	B_\theta(W) = \E_\theta(W) - \theta.
	\]
	Quando $B_\theta(W) = 0$, diz-se que o estimador $W$ é não viesado.
\end{definition}

\begin{definition}[ENVVUM]
	Um estimador $W^*$ é dito ser o melhor estimador não viesado de $\tau(\theta)$ se $\E_\theta(W^*) = \tau(\theta)$ para todo $\theta \in \Theta$ e, para qualquer outro estimador não viesado $W$ de $\tau(\theta)$, $\Var_\theta(W^*) \leq \Var_\theta(W)$ para todo $\theta \in \Theta$. Nesse caso, $W^*$ é chamado de estimador não viesado de variância uniformemente mínima (ENVVUM) de $\tau(\theta)$.
\end{definition}

\begin{theorem}[Desigualdade de Cramér-Rao (Teorema 1)]
	Sejam $X_1, X_2, \ldots, X_n$ uma amostra com f.d.p. conjunta $f(\v{x} | \theta)$ e $W(\v{X})$ um estimador que satisfaz
	\[
	\frac{d}{d\theta} \E_\theta[W(\v{X})] = \int_{\mathcal{X}} \frac{\partial}{\partial\theta} [W(\v{x})f(\v{x} | \theta)]d\v{x} \quad \text{e} \quad \Var_\theta[W(\v{X})] < \infty.
	\]
	Então,
	\[
	\Var_\theta[W(\v{X})] \geq \frac{\left(\frac{d}{d\theta} \E_\theta[W(\v{X})]\right)^2}{\E_\theta\left[\left(\frac{\partial}{\partial\theta} \log f(\v{X} | \theta)\right)^2\right]}.
	\]
	
	\begin{note}
		\leavevmode
		\begin{outline}[enumerate]
			\1 O resultado do Teorema 1, bem como do Corolário 1, permanecerá válido quando as v.a.'s forem discretas. Nesse caso, troca-se a integral pelo somatório na suposição.
	
			\1 A quantidade
			\[
			\frac{\left(\frac{d}{d\theta} \E_\theta[W(\v{X})]\right)^2}{\E_\theta\left[\left(\frac{\partial}{\partial\theta} \log f(\v{X} | \theta)\right)^2\right]},
			\]
			exibida no Teorema 1, é conhecida como limite inferior de Cramér-Rao (LICR).
		
			\1 A quantidade
			\[
			\E_\theta\left[\left(\frac{\partial}{\partial\theta} \log f(\v{X} | \theta)\right)^2\right]
			\]
			é conhecida como informação de Fisher da amostra. Quanto maior a informação de Fisher, mais informação temos sobre $\theta$ e, consequentemente, menor será o limite inferior de Cramér-Rao.
		\end{outline}
	\end{note}
\end{theorem}

\begin{theorem}[Desigualdade de Cramér-Rao -- caso i.i.d. (Corolário 1)]
	Sejam $X_1, X_2, \ldots, X_n$ uma A.A. de uma população com f.d.p. (ou f.p.) $f(x | \theta)$. Sob as suposições do Teorema 1,
	\[
	\Var_\theta[W(\v{X})] \geq \frac{\left(\frac{d}{d\theta} \E_\theta[W(\v{X})]\right)^2}{n\E_\theta\left[\left(\frac{\partial}{\partial\theta} \log f(X | \theta)\right)^2\right]}.
	\]
\end{theorem}

\begin{theorem}[Lema 1]
	Se $f(x | \theta)$ satisfaz
	\[
	\frac{d}{d\theta} \E_\theta\left[\frac{\partial}{\partial\theta} \log f(X | \theta)\right] = \int_{\Re} \frac{\partial}{\partial\theta}\left[\frac{\partial}{\partial\theta} \log f(X | \theta)\right] f(x | \theta) dx,
	\]
	então,
	\[
	\E_\theta\left[\left(\frac{\partial}{\partial\theta} \log f(X | \theta)\right)^2\right] = -\E_\theta\left[\frac{\partial^2}{\partial\theta^2} \log f(X | \theta)\right].
	\]
	
	\begin{note}
		\begin{outline}[enumerate]
			\leavevmode
			\1 A suposição utilizada no Lema 1 é válida, por exemplo, no caso da família exponencial.
			
			\1 Em geral, se o suporte de $X$ depender do parâmetro, a desigualdade de Cramér-Rao não se aplica.
		\end{outline}
	\end{note}
\end{theorem}

\begin{theorem}[Corolário 2]
	Sejam $X_1, X_2, \ldots, X_n$ uma A.A. de uma população $X$ com f.d.p. (ou f.p.) $f(x | \theta)$, em que $f$ satisfaz as condições do Teorema 1. Suponhamos que $L(\theta | \v{x}) = \prod_{i=1}^{n} f(x_i | \theta)$ denote a função de verossimilhança. Se $W(\v{X})$ for um estimador para $\tau(\theta)$, então $W(\v{X})$ atingirá o limite inferior de Cramér-Rao se, e somente se,
	\[
	\frac{\partial}{\partial\theta} \log L(\theta | \v{x}) = a(\theta)[W(\v{x}) - \tau(\theta)],
	\]
	para alguma função $a(\theta)$.
\end{theorem}

\begin{theorem}[Rao-Blackwell (Teorema 2)]
	Seja $W$ qualquer estimador não viesado de $\tau(\theta)$, e $T$ uma estatística suficiente para $\theta$. Defina $\phi(T) = \E[W | T]$. Então $\E_\theta[\phi(T)] = \tau(\theta)$ e $\Var_\theta[\phi(T)] \leq \Var_\theta(W)$ para todo $\theta \in \Theta$. Em outras palavras, $\phi(T)$ é um melhor estimador não viesado uniformemente de $\tau(\theta)$.
	
	\begin{note}
		Devemos considerar estimadores que sejam funções de estatísticas suficientes para determinar ENVVUM's.
	\end{note}
\end{theorem}

\begin{theorem}[Unicidade do melhor ENV (Teorema 3)]
	Se $W$ for um ENVVUM de $\tau(\theta)$, então $W$ será único.
\end{theorem}

\begin{definition}[Estimador não viesado de zero]
	Seja $U$ um estimador tal que $\E_\theta(U) = 0$, para todo $\theta \in \Theta$. Nesse caso, dizemos que $U$ é um estimador não viesado de zero.
\end{definition}

\begin{theorem}[Teorema 4]
	Se $\E_\theta(W) = \tau(\theta)$, então $W$ será o melhor estimador não viesado de $\tau(\theta)$ se, e somente se, $W$ for não correlacionado com todos os estimadores não viesados de zero.
	
	\begin{note}
		Um estimador não viesado de zero nada mais é do que uma perturbação aleatória. Se um estimador pudesse ser melhorado ao acrescentarmos uma perturbação a ele, esse estimador certamente seria problemático. Isso pode ser entendido como uma interpretação do Teorema 4. O Teorema 4 pode ser útil para mostrar quando um estimador não é ENVVUM.
	\end{note}
\end{theorem}

\begin{theorem}[Teorema 5]
	Seja $T$ uma estatística suficiente e completa para um parâmetro $\theta$, e seja $\phi(T)$ qualquer estimador baseado somente em $T$. Então, $\phi(T)$ é o único melhor estimador não viesado de seu valor esperado.
	
	\begin{note}
		O teorema acima diz que, na presença de estatísticas suficientes e completas, se pudermos determinar qualquer estimador não viesado de $\tau(\theta)$, podemos determinar o melhor ENV. Basicamente, se $T$ for uma estatística suficiente e completa para $\theta$ e $h(\v{X})$ um ENV de $\tau(\theta)$, então $\phi(T) = \E[h(\v{X}) | T]$ será o melhor estimador não viesado de $\tau(\theta)$.
	\end{note}
\end{theorem}

\section{Estimação intervalar}
% !TeX root = main.tex
% !TeX spellcheck = pt_BR

% ===========================================================
% Estimação Intervalar
% Fonte: 8 - Estimação Intervalar.pdf
% ===========================================================

% ---------- Definição 1 (Intervalo de confiança) ----------
\begin{definition}
Sejam $X_1,X_2,\ldots,X_n$ uma amostra aleatória de uma população com f.d.p. $f(x\mid\theta)$. Sejam $T_1\equiv T_1(\v{X})$ e $T_2\equiv T_2(\v{X})$ duas estatísticas satisfazendo $T_1\leq T_2$, de tal forma que $\P_\theta[T_1\leq\tau(\theta)\leq T_2]=\gamma$, em que $\gamma$ não depende de $\theta$. Então, o intervalo $[T_1,T_2]$ é chamado de \underline{intervalo com $100\gamma\%$ de confiança para $\tau(\theta)$}. A quantidade $\gamma$ é chamada de \underline{nível de confiança}, e $T_1$ e $T_2$ são chamadas de \underline{limites de confiança inferior e superior}, respectivamente, para $\tau(\theta)$.
\end{definition}
% Correspondência: 8 - Estimação Intervalar.pdf, p. 4/55

% ---------- Definição 2 (IC unilateral) ----------
\begin{definition}
Sejam $X_1,X_2,\ldots,X_n$ uma amostra aleatória de uma população com f.d.p. $f(x\mid\theta)$. Seja $T_1\equiv T_1(\v{X})$ uma estatística tal que $\P_\theta[T_1\leq\tau(\theta)]=\gamma$. Então, $T_1$ é chamada de \underline{limite de confiança inferior unilateral para $\tau(\theta)$}. Analogamente, seja $T_2\equiv T_2(\v{X})$ uma estatística tal que $\P_\theta[T_2\geq\tau(\theta)]=\gamma$. Então, $T_2$ é chamada de \underline{limite de confiança superior unilateral para $\tau(\theta)$}. Em ambos os casos, $\gamma$ não depende de $\theta$.
\end{definition}
% Correspondência: 8 - Estimação Intervalar.pdf, p. 5/55

% ---------- Observação 1 ----------
\begin{note}
Se um IC para $\theta$ tiver sido determinado, então pode-se determinar uma família de IC's. Em outras palavras, para um determinado estimador intervalar com $100\gamma\%$ de confiança para $\theta$, pode-se obter um estimador intervalar com $100\gamma\%$ para $\tau(\theta)$, desde que $\tau$ seja uma função estritamente monótona. Se $\tau$ for uma função estritamente crescente e $[T_1,T_2]$ for um IC para $\theta$, então $[\tau(T_1),\tau(T_2)]$ será um IC para $\tau(\theta)$, pois
\begin{equation}
\P_\theta[T_1\leq\theta\leq T_2] = \P_\theta[\tau(T_1)\leq\tau(\theta)\leq\tau(T_2)].
\end{equation}
\end{note}
% Correspondência: 8 - Estimação Intervalar.pdf, p. 7/55

% ---------- Definição 3 (Quantidade pivotal) ----------
\begin{definition}
Sejam $X_1,X_2,\ldots,X_n$ uma A.A. de uma população com f.d.p. $f(x\mid\theta)$, e seja $Q\equiv Q(\v{X}\mid\theta)$. Dizemos que $Q$ é uma \underline{quantidade pivotal} se sua distribuição não depender de $\theta$.
\end{definition}
% Correspondência: 8 - Estimação Intervalar.pdf, p. 10/55

% ---------- Método da quantidade pivotal ----------
\begin{theorem}[Método da quantidade pivotal]
Se $Q(\v{X}\mid\theta)$ é uma quantidade pivotal e possui uma f.d.p., então, para qualquer $0<\gamma<1$, haverá $q_1$ e $q_2$ dependendo de $\gamma$ de tal modo que
\begin{equation}
\P(q_1\leq Q(\v{X}\mid\theta)\leq q_2) = \gamma.
\end{equation}
Se para cada possível valor amostral $x_1,x_2,\ldots,x_n$,
\begin{equation}
\{q_1\leq Q(\v{x}\mid\theta)\leq q_2\} \iff \{T_1(\v{x})\leq\tau(\theta)\leq T_2(\v{x})\},
\end{equation}
então $[T_1(\v{X}),T_2(\v{X})]$ será um intervalo com $100\gamma\%$ de confiança para $\tau(\theta)$.
\end{theorem}
% Correspondência: 8 - Estimação Intervalar.pdf, p. 12/55

% ---------- Observação 2 e 3 ----------
\begin{note}
\begin{outline}
	\1 O termo quantidade pivotal pode ser inapropriado, pois pode ser impossível ``pivotar'' a quantidade $Q(\v{X}\mid\theta)$.
	\1 Os valores de $q_1$ e $q_2$ não dependem de $\theta$, pois a distribuição de $Q(\v{X}\mid\theta)$ não depende.
\end{outline}
\end{note}
% Correspondência: 8 - Estimação Intervalar.pdf, p. 13/55

% ---------- Utilidade pública: relações (1) e (2) ----------
\begin{note}
No caso do método da quantidade pivotal, intervalos com $100\gamma\%$ de confiança são caracterizados por
\begin{equation}
\gamma = \P(q_1\leq Q\leq q_2) = \int_{q_1}^{q_2} f_Q(t)\,dt = F_Q(q_2)-F_Q(q_1), \label{eq:util-pub-1}
\end{equation}
em que $f_Q$ e $F_Q$ correspondem à f.d.p. e f.d. da quantidade pivotal $Q$. Disso decorrem duas informações úteis:
\begin{outline}
	\1 Como $\gamma$ é fixado, podemos escrever $q_2\equiv q_2(q_1)$, pois
	\begin{equation}
	\gamma = F_Q(q_2)-F_Q(q_1) \implies q_2 = F_Q^{-1}[\gamma+F_Q(q_1)];
	\end{equation}
	\1 Derivando ambos os lados de \eqref{eq:util-pub-1} em relação a $q_1$, temos
	\begin{equation}
	\frac{d}{dq_1}\gamma = \frac{d}{dq_1}\left[F_Q(q_2)-F_Q(q_1)\right] \implies f_Q(q_2)\frac{dq_2}{dq_1}-f_Q(q_1) = 0.
	\end{equation}
	Portanto,
	\begin{equation}
	\frac{dq_2}{dq_1} = \frac{f_Q(q_1)}{f_Q(q_2)}. \label{eq:dq2dq1}
	\end{equation}
\end{outline}
\end{note}
% Correspondência: 8 - Estimação Intervalar.pdf, p. 17-18/55

% ---------- Quantidade pivotal para a média, sigma^2 conhecido ----------
\begin{theorem}
Sejam $X_1,X_2,\ldots,X_n$ uma A.A. de uma população $X\sim\normdist{\theta,\sigma^2}$, $\theta\in\Re$ e $\sigma^2>0$ conhecido. Então
\begin{equation}
Q = \frac{\bar{X}-\theta}{\sigma/\sqrt{n}} \sim \normdist{0,1}
\end{equation}
é uma quantidade pivotal e
\begin{equation}
\left[\bar{X}-q_2\frac{\sigma}{\sqrt{n}},\ \bar{X}-q_1\frac{\sigma}{\sqrt{n}}\right]
\end{equation}
é um IC de $100\gamma\%$ para $\theta$, em que $q_1$ e $q_2$ satisfazem $\P(q_1\leq Q\leq q_2)=\gamma$. O comprimento desse intervalo é minimizado quando $q_1=-q_2$.
\end{theorem}
% Correspondência: 8 - Estimação Intervalar.pdf, p. 11 e 15/55

% ---------- IC para a média, sigma^2 desconhecido ----------
\begin{theorem}[IC para a média — $\sigma^2$ desconhecido]
Sejam $X_1,X_2,\ldots,X_n$ uma A.A. oriunda de uma população $X\sim\normdist{\mu,\sigma^2}$, com $\sigma^2$ desconhecido. Então
\begin{equation}
Q = \frac{\bar{X}-\mu}{S/\sqrt{n}} \sim t_{n-1}
\end{equation}
é uma quantidade pivotal, e o menor intervalo com $100\gamma\%$ de confiança para $\mu$ é
\begin{equation}
\left[\bar{X}-t_{n-1,(1+\gamma)/2}\frac{S}{\sqrt{n}},\ \bar{X}+t_{n-1,(1+\gamma)/2}\frac{S}{\sqrt{n}}\right],
\end{equation}
em que $t_{n-1,(1+\gamma)/2}$ é tal que $\P(T\leq t_{n-1,(1+\gamma)/2})=(1+\gamma)/2$, para $T\sim t_{n-1}$.
\end{theorem}
% Correspondência: 8 - Estimação Intervalar.pdf, p. 19-21 e 25/55

% ---------- IC para a variância, mu desconhecido ----------
\begin{theorem}[IC para a variância — $\mu$ desconhecido]
Sejam $X_1,X_2,\ldots,X_n$ uma A.A. oriunda de uma população $X\sim\normdist{\mu,\sigma^2}$, com $\mu$ desconhecido. Então
\begin{equation}
Q = \frac{(n-1)S^2}{\sigma^2} \sim \chidist{n-1}
\end{equation}
é uma quantidade pivotal, e
\begin{equation}
\left[\frac{(n-1)S^2}{q_2},\ \frac{(n-1)S^2}{q_1}\right]
\end{equation}
é um intervalo com $100\gamma\%$ de confiança para $\sigma^2$, em que $q_1$ e $q_2$ satisfazem $\P(q_1\leq Q\leq q_2)=\gamma$.
\begin{outline}
	\1 Uma escolha usual (\underline{intervalo com caudas iguais}) é tomar $q_1$ e $q_2$ tais que $\P(Q\leq q_1)=\P(Q\geq q_2)=(1-\gamma)/2$, mas esse não é o intervalo de comprimento mínimo.
	\1 O intervalo de comprimento mínimo é obtido sujeitando-se $q_1$ e $q_2$ à condição adicional $q_1^2f_Q(q_1)=q_2^2f_Q(q_2)$, cuja solução é obtida numericamente.
\end{outline}
\end{theorem}
% Correspondência: 8 - Estimação Intervalar.pdf, p. 26-31/55

% ---------- IC para a diferença de médias, variâncias iguais e desconhecidas ----------
\begin{theorem}[IC para a diferença de médias — $\sigma_1^2=\sigma_2^2=\sigma^2$ desconhecido]
Sejam $X_{11},X_{12},\ldots,X_{1n_1}$ uma A.A. de uma população $X_1\sim\normdist{\mu_1,\sigma_1^2}$ e sejam $X_{21},X_{22},\ldots,X_{2n_2}$ uma A.A. de uma população $X_2\sim\normdist{\mu_2,\sigma_2^2}$, ambas as amostras independentes, com $\sigma_1^2=\sigma_2^2=\sigma^2$ desconhecido. Sejam
\begin{equation}
S_p^2 = \frac{(n_1-1)S_1^2+(n_2-1)S_2^2}{n_1+n_2-2}.
\end{equation}
Então
\begin{equation}
Q = \frac{(\bar{X}_1-\bar{X}_2)-(\mu_1-\mu_2)}{\sqrt{\left(\dfrac{1}{n_1}+\dfrac{1}{n_2}\right)S_p^2}} \sim t_{n_1+n_2-2}
\end{equation}
é uma quantidade pivotal, e
\begin{equation}
\left[(\bar{X}_1-\bar{X}_2)-q\sqrt{\left(\frac{1}{n_1}+\frac{1}{n_2}\right)S_p^2},\ (\bar{X}_1-\bar{X}_2)+q\sqrt{\left(\frac{1}{n_1}+\frac{1}{n_2}\right)S_p^2}\right]
\end{equation}
é um intervalo com $100\gamma\%$ de confiança para $\mu_1-\mu_2$, em que $q=t_{n_1+n_2-2,(1+\gamma)/2}$.
\end{theorem}
% Correspondência: 8 - Estimação Intervalar.pdf, p. 32-37/55

% ---------- IC para a diferença de médias, amostras pareadas ----------
\begin{theorem}[IC para a diferença de médias — amostras pareadas, $\sigma_D^2$ desconhecida]
Sejam $(X_{11},X_{12}),(X_{21},X_{22}),\ldots,(X_{n1},X_{n2})$ uma A.A. oriunda de uma população normal bivariada com parâmetros $\E(X_1)=\mu_1$, $\E(X_2)=\mu_2$, $\Var(X_1)=\sigma_1^2$, $\Var(X_2)=\sigma_2^2$ e $\rho=\operatorname{Cor}(X_1,X_2)$. Sejam $D_i=X_{i1}-X_{i2}$, $i=1,2,\ldots,n$, com $D_i\sim\normdist{\mu_D,\sigma_D^2}$, em que $\mu_D=\mu_1-\mu_2$ e $\sigma_D^2=\sigma_1^2+\sigma_2^2-2\operatorname{Cov}(X_{i1},X_{i2})$. Então
\begin{equation}
Q = \frac{\bar{D}-\mu_D}{S_D/\sqrt{n}} \sim t_{n-1}
\end{equation}
é uma quantidade pivotal, e
\begin{equation}
\left[\bar{D}-t_{n-1,(1+\gamma)/2}\frac{S_D}{\sqrt{n}},\ \bar{D}+t_{n-1,(1+\gamma)/2}\frac{S_D}{\sqrt{n}}\right]
\end{equation}
é um intervalo com $100\gamma\%$ de confiança para $\mu_D=\mu_1-\mu_2$.
\end{theorem}
% Correspondência: 8 - Estimação Intervalar.pdf, p. 38-40/55

% ---------- Existência de quantidade pivotal (caso geral) ----------
\begin{theorem}[Existência de quantidade pivotal]
Sejam $X_1,X_2,\ldots,X_n$ uma A.A. oriunda de uma população com f.d.p. $f(x\mid\theta)$ e f.d.a. contínua $F(x\mid\theta)$. Então
\begin{equation}
Q = -\sum_{i=1}^n \log F(X_i\mid\theta) \sim \operatorname{gama}(n,1)
\end{equation}
é uma quantidade pivotal, equivalente a
\begin{equation}
Q' = \prod_{i=1}^n F(X_i\mid\theta).
\end{equation}
Ou seja, sempre que a amostra vier de uma população com f.d.a. contínua, há uma quantidade pivotal.
\end{theorem}
% Correspondência: 8 - Estimação Intervalar.pdf, p. 42-45/55

% ---------- Observações 4 e 5 (caso geral) ----------
\begin{note}
\begin{outline}
	\1 Embora seja possível determinar uma quantidade pivotal sempre que uma amostra for oriunda de uma população com f.d.a. contínua, não há garantias de que seja possível determinar um IC a partir dela. Em outras palavras, pode ser impossível ``pivotar'' essa quantidade.
	\1 Se $F(x\mid\theta)$ for uma função monótona em $\theta$ para cada $x$ fixado, então $\prod_{i=1}^n F(X_i\mid\theta)$ também será. Essa é uma condição suficiente para que seja possível determinar intervalos de confiança.
\end{outline}
\end{note}
% Correspondência: 8 - Estimação Intervalar.pdf, p. 46/55

% ---------- Famílias de locação-escala ----------
\begin{note}
Se $\theta$ for um parâmetro de locação, então a distribuição de $X_i-\theta$ não depende de $\theta$; exemplos de quantidade pivotal são $\sum_{i=1}^n X_i-n\theta$, $\bar{X}-\theta$, $X_{(k)}-\theta$, $\dfrac{X_{(1)}+X_{(n)}}{2}-\theta$, etc.

Se $\theta$ for um parâmetro de escala, a distribuição de $X_i/\theta$ não depende de $\theta$; exemplos de quantidade pivotal são $\sum_{i=1}^n X_i/\theta$, $X_{(k)}/\theta$, etc.
\end{note}
% Correspondência: 8 - Estimação Intervalar.pdf, p. 48/55

% ---------- Distribuição assintótica de EMV's ----------
\begin{theorem}[Distribuição assintótica de EMV's]
Sejam $X_1,X_2,\ldots,X_n$ uma A.A. oriunda de uma população com f.d.a. $F(x\mid\theta)$. Denotemos por $\hat{\theta}_n$ o EMV de $\theta$. Sob certas condições de regularidade,
\begin{equation}
\frac{\hat{\theta}_n-\theta}{\sqrt{\sigma_n^2(\theta)}} \toD \normdist{0,1},
\end{equation}
quando $n\to\infty$, em que $1/\sigma_n^2(\theta)=I_n(\theta)$ corresponde à informação de Fisher, ou seja,
\begin{equation}
I_n(\theta) = n\E_\theta\left\{\left[\frac{\partial}{\partial\theta}\log f(X\mid\theta)\right]^2\right\} = -n\E_\theta\left[\frac{\partial^2}{\partial\theta^2}\log f(X\mid\theta)\right].
\end{equation}
\end{theorem}
% Correspondência: 8 - Estimação Intervalar.pdf, p. 51/55

% ---------- IC's assintóticos para grandes amostras ----------
\begin{note}
Seja $\hat{\theta}_n$ um estimador tal que
\begin{equation}
Q_n = \frac{\hat{\theta}_n-\theta}{\sigma_n(\theta)} \toD \normdist{0,1},
\end{equation}
quando $n\to\infty$, em que $\sigma_n(\theta)$ é uma função de $\theta$ que depende de $n$. Então $Q_n$ pode ser tratada como uma quantidade pivotal aproximada, e um IC com um nível de aproximadamente $100\gamma\%$ de confiança pode ser determinado invertendo
\begin{equation}
-z_{(1+\gamma)/2} \leq Q_n \leq z_{(1+\gamma)/2},
\end{equation}
em que $\P(Z\leq z_{(1+\gamma)/2})=(1+\gamma)/2$, desde que a quantidade pivotal $Q_n$ possa ser ``pivotada'' na desigualdade acima.
\end{note}
% Correspondência: 8 - Estimação Intervalar.pdf, p. 52-53/55


\section{Testes de hipóteses I}
% !TeX root = main.tex
% !TeX spellcheck = pt_BR
% ===========================================================
% Testes de Hipóteses I: conceitos iniciais
% Fonte: 9 - Testes de Hipóteses 1.pdf
% ===========================================================

% ---------- Definição 1 (Hipótese estatística) ----------
\begin{definition}
Uma \uline{hipótese estatística} é uma afirmação ou conjectura a respeito da distribuição de uma ou mais variáveis. Se uma hipótese estatística especifica completamente a distribuição dos dados, então ela é chamada de \uline{hipótese simples}; caso contrário, ela é chamada de \uline{hipótese composta}.
\end{definition}
% Correspondência: 9 - Testes de Hipóteses 1.pdf, p. 4/45

% ---------- Exemplo 1 (População normal – variância conhecida) ----------
\begin{example}[População normal – variância conhecida]
Sejam $X_1,X_2,\ldots,X_n$ uma A.A. de uma população $X\sim\normdist{\theta,25}$. A hipótese estatística de que a média dessa população normal é menor do que ou igual a 17 é denotada por
\begin{equation}
\mathcal{H}:\theta\leq17.
\end{equation}
Essa hipótese é composta, porque não especifica completamente a distribuição dos dados. Por outro lado, a hipótese
\begin{equation}
\mathcal{H}:\theta=17
\end{equation}
é simples, já que especifica completamente a distribuição populacional.
\end{example}
% Correspondência: 9 - Testes de Hipóteses 1.pdf, p. 5/45

% ---------- Definição 2 (Teste) ----------
\begin{definition}
Um \uline{teste} $\Upsilon$ de uma hipótese estatística $\mathcal{H}$ é uma regra, ou procedimento, para decidir se $\mathcal{H}$ deve ser rejeitada.
\end{definition}
% Correspondência: 9 - Testes de Hipóteses 1.pdf, p. 6/45

% ---------- Exemplo 2 (Cont. Ex. 1 – dois possíveis testes) ----------
\begin{example}[Cont. Exemplo 1]
Sejam $X_1,X_2,\ldots,X_n$ uma A.A. de uma população $X\sim\normdist{\theta,25}$. Consideremos $\mathcal{H}:\theta\leq17$.

Um possível teste, $\Upsilon$, seria \uline{rejeitar $\mathcal{H}$ se, e somente se,} $\bar{x}>17+5/\sqrt{n}$.

Outro possível teste, $\Upsilon'$, seria \uline{lançar uma moeda e rejeitar $\mathcal{H}$ se, e somente se, o resultado for cara}.
\end{example}
% Correspondência: 9 - Testes de Hipóteses 1.pdf, p. 7/45

% ---------- Observação 1 (Espaço amostral) ----------
\begin{note}
Lembrete: utilizamos $\mathcal{X}$ para denotar o \uline{espaço amostral} dos dados, ou seja,
\begin{equation}
\mathcal{X}=\{(x_1,x_2,\ldots,x_n)\in\mathbb{R}^n\mid f(x_1,x_2,\ldots,x_n\mid\theta)>0\}=\{\v{x}\in\mathbb{R}^n\mid f(\v{x}\mid\theta)>0\}.
\end{equation}
\end{note}
% Correspondência: 9 - Testes de Hipóteses 1.pdf, p. 8/45

% ---------- Definição 3 (Teste não aleatorizado) ----------
\begin{definition}
Seja $\Upsilon$ o teste de uma hipótese $\mathcal{H}$ definido como \uline{rejeitar $\mathcal{H}$ se, e somente se,} $(x_1,x_2,\ldots,x_n)\in\mathcal{C}_\Upsilon$, em que $\mathcal{C}_\Upsilon$ corresponde a um subconjunto de $\mathcal{X}$. Então, $\Upsilon$ é chamado de \uline{teste não aleatorizado} e $\mathcal{C}_\Upsilon$ é denominada \uline{região crítica} do teste $\Upsilon$.
\end{definition}
% Correspondência: 9 - Testes de Hipóteses 1.pdf, p. 9/45

% ---------- Exemplo 3 (Continuação – população normal) ----------
\begin{example}[Continuação – população normal]
Sejam $X_1,X_2,\ldots,X_n$ uma A.A. de uma população $X\sim\normdist{\theta,25}$. Nesse caso, $\mathcal{X}=\mathbb{R}^n$.

Consideremos $\mathcal{H}:\theta\leq17$ e o teste $\Upsilon$: \uline{rejeitar $\mathcal{H}$ se, e somente se,}
\begin{equation}
\bar{x}>17+5/\sqrt{n}.
\end{equation}
Nesse caso, $\Upsilon$ é não aleatorizado, e
\begin{equation}
\mathcal{C}_\Upsilon=\{\v{x}\in\mathcal{X}\mid\bar{x}>17+5/\sqrt{n}\}.
\end{equation}
\end{example}
% Correspondência: 9 - Testes de Hipóteses 1.pdf, p. 10/45

% ---------- Observação 2 (Decomposição do espaço amostral) ----------
\begin{note}
Um teste não aleatorizado $\Upsilon$ de $\mathcal{H}$ consiste em uma decomposição
\begin{equation}
\mathcal{X}=\mathcal{C}_\Upsilon\cup\bar{\mathcal{C}}_\Upsilon,
\end{equation}
com $\mathcal{C}_\Upsilon\cap\bar{\mathcal{C}}_\Upsilon=\emptyset$, de modo que, se uma amostra observada $\v{x}\in\mathcal{C}_\Upsilon$, então rejeita-se $\mathcal{H}$. Assim, um teste não aleatorizado é completamente especificado por sua região crítica correspondente.
\end{note}
% Correspondência: 9 - Testes de Hipóteses 1.pdf, p. 11/45

% ---------- Definição 4 (Teste aleatorizado) ----------
\begin{definition}
Um teste $\Upsilon$ de uma hipótese $\mathcal{H}$ é chamado de \uline{teste aleatorizado} se $\Upsilon$ for definido pela função
\begin{equation}
\Psi_\Upsilon(\v{x})=\P(\mathcal{H}\text{ ser rejeitada}\mid\v{x}\text{ é observada}).
\end{equation}
A função $\Psi_\Upsilon$ é chamada de \uline{função crítica} do teste $\Upsilon$.
\end{definition}
% Correspondência: 9 - Testes de Hipóteses 1.pdf, p. 12/45

% ---------- Exemplo 4 (População Bernoulli) ----------
\begin{example}[População Bernoulli]
Sejam $X_1,X_2,\ldots,X_{10}$ uma A.A. de uma população $X\sim\berndist{\theta}$. Suponha que estejamos interessados em testar a hipótese $\mathcal{H}:\theta<1/2$.

Um possível teste $\Upsilon$ seria:
\begin{outline}
\1 aceitar $\mathcal{H}$ se $\sum_{i=1}^{10}x_i<5$;
\1 rejeitar $\mathcal{H}$ se $\sum_{i=1}^{10}x_i>5$;
\1 decidir entre rejeitar e aceitar $\mathcal{H}$ com o lançamento de uma moeda, se $\sum_{i=1}^{10}x_i=5$.
\end{outline}
O lançamento da moeda é o ensaio auxiliar de Bernoulli.

Esse teste $\Upsilon$ particiona $\mathcal{X}$ em três regiões:
\begin{align}
A &= \left\{(x_1,x_2,\ldots,x_{10})\in\mathcal{X}\,\middle|\,\sum_{i=1}^{10}x_i<5\right\},\\
B &= \left\{(x_1,x_2,\ldots,x_{10})\in\mathcal{X}\,\middle|\,\sum_{i=1}^{10}x_i=5\right\},\\
C &= \left\{(x_1,x_2,\ldots,x_{10})\in\mathcal{X}\,\middle|\,\sum_{i=1}^{10}x_i>5\right\}.
\end{align}

A função crítica do teste $\Upsilon$ é dada por
\begin{equation}
\Psi_\Upsilon(x_1,x_2,\ldots,x_{10})=
\begin{cases}
1, & \text{se }(x_1,x_2,\ldots,x_{10})\in C;\\
1/2, & \text{se }(x_1,x_2,\ldots,x_{10})\in B;\\
0, & \text{se }(x_1,x_2,\ldots,x_{10})\in A.
\end{cases}
\end{equation}
\end{example}
% Correspondência: 9 - Testes de Hipóteses 1.pdf, p. 15/45

% ---------- Definição 5 (Erros do tipo I e II) ----------
\begin{definition}
A rejeição de $\mathcal{H}_0$ quando ela é verdadeira é chamada de \uline{erro do tipo I}, e a aceitação de $\mathcal{H}_0$ quando ela é falsa é chamada de \uline{erro do tipo II}. O \uline{tamanho de um erro do tipo I} é definido como a probabilidade de se cometer um erro do tipo I. Similarmente, o \uline{tamanho de um erro do tipo II} é a probabilidade de se cometer um erro do tipo II.
\end{definition}
% Correspondência: 9 - Testes de Hipóteses 1.pdf, p. 19/45

% ---------- Definição 6 (Função poder) ----------
\begin{definition}
Seja $\Upsilon$ um teste com hipótese nula $\mathcal{H}_0$. A \uline{função poder} do teste $\Upsilon$, denotada por $\pi_\Upsilon(\theta)$, é definida como sendo a probabilidade de que $\mathcal{H}_0$ seja rejeitada quando a distribuição da qual a amostra foi selecionada tem parâmetro $\theta$.
\end{definition}
% Correspondência: 9 - Testes de Hipóteses 1.pdf, p. 20/45

% ---------- Exemplo 5 (Continuação – população normal, função poder) ----------
\begin{example}[Continuação – população normal]
Sejam $X_1,X_2,\ldots,X_n$ uma A.A. de uma população $X\sim\normdist{\theta,25}$.

Consideremos $\mathcal{H}_0:\theta\leq17$ e o teste $\Upsilon$: rejeitar $\mathcal{H}_0$ se, e somente se, $\bar{x}>17+5/\sqrt{n}$.

Nesse caso,
\begin{equation}
\pi_\Upsilon(\theta)=\P_\theta\left(\bar{X}>17+5/\sqrt{n}\right)=1-\Phi\left(\frac{17+5/\sqrt{n}-\theta}{5/\sqrt{n}}\right).
\end{equation}
% Figura: gráfico da função poder $\pi_\Upsilon(\theta)$ para $n=25$ (crescente em $\theta$, de $0$ a $1$, com $\pi_\Upsilon(17)\approx0{,}16$ e $\pi_\Upsilon(20)\approx1$).
\end{example}
% Correspondência: 9 - Testes de Hipóteses 1.pdf, p. 23/45

% ---------- Observação 3 (Função poder para diferentes testes) ----------
\begin{note}
Pode-se escrever
\begin{equation}
\pi_\Upsilon(\theta)=\P_\theta(\text{rejeitar }\mathcal{H}_0),
\end{equation}
em que $\theta$ é o verdadeiro parâmetro.

Se $\Upsilon$ for um teste não aleatorizado, então
\begin{equation}
\pi_\Upsilon(\theta)=\P_\theta(\v{X}\in\mathcal{C}_\Upsilon),
\end{equation}
em que $\mathcal{C}_\Upsilon$ corresponde à região crítica associada ao teste $\Upsilon$.

Caso $\Upsilon$ seja um teste aleatorizado com função crítica $\Psi_\Upsilon$, então
\begin{align}
\pi_\Upsilon(\theta) &= \P_\theta(\text{rejeitar }\mathcal{H}_0)\notag\\
&= \int_{\mathcal{X}} \P_\theta(\text{rejeitar }\mathcal{H}_0\mid\v{x})f_{\v{X}}(\v{x}\mid\theta)\,d\v{x}\notag\\
&= \int_{\mathcal{X}} \Psi_\Upsilon(\v{x})f_{\v{X}}(\v{x}\mid\theta)\,d\v{x}\notag\\
&= \E_\theta\left[\Psi_\Upsilon(\v{X})\right].
\end{align}
\end{note}
% Correspondência: 9 - Testes de Hipóteses 1.pdf, p. 22/45

% ---------- Definição 7 (Tamanho do teste / nível de significância) ----------
\begin{definition}
Seja $\Upsilon$ um teste da hipótese $\mathcal{H}_0:\theta\in\Theta_0$, em que $\Theta_0\subset\Theta$. O \uline{tamanho do teste} $\Upsilon$ de $\mathcal{H}_0$ é definido como
\begin{equation}
\alpha=\sup_{\theta\in\Theta_0}\pi_\Upsilon(\theta).
\end{equation}
O tamanho do teste no caso não aleatorizado também é conhecido como \uline{tamanho da região crítica} e, comumente, como \uline{nível de significância}.
\end{definition}
% Correspondência: 9 - Testes de Hipóteses 1.pdf, p. 24/45

% ---------- Exemplo 6 (Continuação – população normal, tamanho do teste) ----------
\begin{example}[Continuação – população normal]
Sejam $X_1,X_2,\ldots,X_n$ uma A.A. de uma população $X\sim\normdist{\theta,25}$.

Consideremos $\mathcal{H}_0:\theta\leq17$ e o teste $\Upsilon$: rejeitar $\mathcal{H}_0$ se, e somente se, $\bar{x}>17+5/\sqrt{n}$.

Nesse caso, $\Theta_0=\{\theta\in\Theta\mid\theta\leq17\}$ e o tamanho do teste $\Upsilon$ será
\begin{equation}
\sup_{\theta\in\Theta_0}\pi_\Upsilon(\theta)=\sup_{\theta\leq17}\left[1-\Phi\left(\frac{17+5/\sqrt{n}-\theta}{5/\sqrt{n}}\right)\right]=1-\Phi(1)\approx0{,}159.
\end{equation}
\end{example}
% Correspondência: 9 - Testes de Hipóteses 1.pdf, p. 25/45

% ---------- Teorema 1 (Suficiência) ----------
\begin{theorem}
Sejam $X_1,X_2,\ldots,X_n$ uma A.A. de uma população com f.d.p. $f(x\mid\theta)$, $\theta\in\Theta$, e $(T_1,T_2,\ldots,T_s)=(T_1(\v{X}),T_2(\v{X}),\ldots,T_s(\v{X}))$ uma estatística suficiente para $\theta$. Então, para qualquer teste $\Upsilon$ com função crítica $\Psi_\Upsilon$, existe um teste, digamos $\Upsilon'$, e uma função crítica correspondente, digamos $\Psi_{\Upsilon'}$, dependendo somente do conjunto de estatísticas suficientes, que satisfazem $\pi_\Upsilon(\theta)=\pi_{\Upsilon'}(\theta)$ para todo $\theta\in\Theta$.
\end{theorem}
\begin{proof}
Definamos
\begin{equation}
\Psi_{\Upsilon'}(t_1,t_2,\ldots,t_s)=\E\left[\Psi_\Upsilon(X_1,X_2,\ldots,X_n)\mid T_1=t_1,T_2=t_2,\ldots,T_s=t_s\right].
\end{equation}
Portanto, $\Psi_{\Upsilon'}$ é uma função crítica. Além disso,
\begin{align}
\pi_{\Upsilon'}(\theta) &= \E_\theta\left[\Psi_{\Upsilon'}(T_1,T_2,\ldots,T_s)\right]\notag\\
&= \E_\theta\left\{\E\left[\Psi_\Upsilon(X_1,X_2,\ldots,X_n)\mid T_1,T_2,\ldots,T_s\right]\right\}\notag\\
&= \pi_\Upsilon(\theta).
\end{align}
\end{proof}
% Correspondência: 9 - Testes de Hipóteses 1.pdf, p. 26-27/45

% ---------- Definição 8 (Teste da razão de verossimilhanças simples) ----------
\begin{definition}
Sejam $X_1,X_2,\ldots,X_n$ uma A.A. de uma população com f.d.p. $f_0$ ou $f_1$. Um teste $\Upsilon$ de $\mathcal{H}_0:X_i\sim f_0$ \emph{versus} $\mathcal{H}_1:X_i\sim f_1$ é chamado de \uline{teste da razão de verossimilhanças simples} se $\Upsilon$ for caracterizado como
\begin{outline}
\1 rejeitar $\mathcal{H}_0$, se $\lambda(\v{x})<k$;
\1 aceitar $\mathcal{H}_0$, se $\lambda(\v{x})>k$;
\1 aceitar ou rejeitar $\mathcal{H}_0$, ou aleatorizar, se $\lambda(\v{x})=k$,
\end{outline}
em que $k$ é uma constante positiva e
\begin{equation}
\lambda(\v{x})=\lambda(x_1,x_2,\ldots,x_n)=\frac{L(\theta_0\mid\v{x})}{L(\theta_1\mid\v{x})}\equiv\frac{L_0}{L_1}.
\end{equation}
\end{definition}
% Correspondência: 9 - Testes de Hipóteses 1.pdf, p. 32/45

% ---------- Definição 9 (Teste mais poderoso) ----------
\begin{definition}
Um teste $\Upsilon^*$ de $\mathcal{H}_0:\theta=\theta_0$ \emph{versus} $\mathcal{H}_1:\theta=\theta_1$ é definido ser um \uline{teste mais poderoso} (TMP) de tamanho $\alpha$ ($0<\alpha<1$) se, e somente se,
\begin{outline}
\1 (i) $\pi_{\Upsilon^*}(\theta_0)=\alpha$;
\1 (ii) $\pi_{\Upsilon^*}(\theta_1)\geq\pi_\Upsilon(\theta_1)$ para qualquer outro teste $\Upsilon$ tal que $\pi_\Upsilon(\theta_0)\leq\alpha$.
\end{outline}
\end{definition}
% Correspondência: 9 - Testes de Hipóteses 1.pdf, p. 37/45

% ---------- Observação 4 (Interpretação de TMP) ----------
\begin{note}
Um teste $\Upsilon^*$ é mais poderoso de tamanho $\alpha$ se o tamanho do seu erro do tipo I for $\alpha$ e o tamanho do seu erro do tipo II for o menor, dentre todos os testes de tamanho $\alpha$ ou menos.
\end{note}
% Correspondência: 9 - Testes de Hipóteses 1.pdf, p. 37/45

% ---------- Teorema 2 (Lema de Neyman-Pearson) ----------
\begin{theorem}[Lema de Neyman-Pearson]
Sejam $X_1,X_2,\ldots,X_n$ uma A.A. de uma população com f.d.p. $f(x\mid\theta)$, $\theta\in\{\theta_0,\theta_1\}$, e seja $0<\alpha<1$ fixado. Denotemos por $k^*$ uma constante positiva e $\mathcal{C}^*$ um subconjunto de $\mathcal{X}$ satisfazendo:
\begin{outline}
\1 (i) $\P_{\theta_0}(\v{X}\in\mathcal{C}^*)=\alpha$;
\1 (ii) $\lambda(\v{x})=\dfrac{L(\theta_0\mid\v{x})}{L(\theta_1\mid\v{x})}=\dfrac{L_0}{L_1}\leq k^*$, se $\v{x}\in\mathcal{C}^*$, e $\lambda(\v{x})\geq k^*$, se $\v{x}\in\bar{\mathcal{C}}^*$, em que $\bar{\mathcal{C}}^*=\mathcal{X}-\mathcal{C}^*$.
\end{outline}
Então, o teste $\Upsilon^*$ com região crítica $\mathcal{C}^*$ é um teste mais poderoso de tamanho $\alpha$ para testar $\mathcal{H}_0:\theta=\theta_0$ \emph{vs.} $\mathcal{H}_1:\theta=\theta_1$.
\end{theorem}
% Correspondência: 9 - Testes de Hipóteses 1.pdf, p. 38/45

% ---------- Observação 5 (Existência de k* e C*) ----------
\begin{note}
$k^*$ e $\mathcal{C}^*$ satisfazendo (i) e (ii) no Teorema 2 nem sempre existem e, consequentemente, o teorema pode não fornecer um TMP de tamanho $\alpha$, pelo menos da forma como foi estabelecido.

Quando a A.A. for composta por v.a.'s contínuas, $k^*$ e $\mathcal{C}^*$ existirão. O problema é quando forem v.a.'s discretas.

Embora o Lema de Neyman-Pearson tenha sido definido apenas para testes não aleatorizados, ele pode ser generalizado, de forma a contemplar também testes aleatorizados.
\end{note}
% Correspondência: 9 - Testes de Hipóteses 1.pdf, p. 39/45

% ---------- Exemplo 7 (Teste mais poderoso – exponencial) ----------
\begin{example}
Sejam $X_1,X_2,\ldots,X_n$ uma A.A. oriunda de uma população $X\sim\text{exp}(\theta)$, em que $\theta\in\{\theta_0,\theta_1\}$, com $\theta_0$ e $\theta_1$ números fixados.

Suponhamos aqui que queremos testar $\mathcal{H}_0:\theta=\theta_0$ \emph{vs.} $\mathcal{H}_1:\theta=\theta_1$, e que $\theta_0<\theta_1$.

Observe que
\begin{equation}
L_k=L(\theta_k\mid\v{x})=\theta_k^n e^{-\theta_k\sum_{i=1}^n x_i}, \quad k=0,1.
\end{equation}
Nesse problema, portanto,
\begin{align}
\lambda(\v{x}) &= \frac{L_0}{L_1}=\frac{\theta_0^n e^{-\theta_0\sum_{i=1}^n x_i}}{\theta_1^n e^{-\theta_1\sum_{i=1}^n x_i}}\notag\\
&= \left(\frac{\theta_0}{\theta_1}\right)^n\exp\left[(\theta_1-\theta_0)\sum_{i=1}^n x_i\right].
\end{align}
Pelo Lema de Neyman-Pearson, o teste mais poderoso rejeita $\mathcal{H}_0$ se, e somente se, $\lambda(\v{x})\leq k^*$.

Assim,
\begin{align}
\left(\frac{\theta_0}{\theta_1}\right)^n\exp\left[(\theta_1-\theta_0)\sum_{i=1}^n x_i\right]\leq k^* &\iff\notag\\
\log\left(\frac{\theta_0}{\theta_1}\right)^n+(\theta_1-\theta_0)\sum_{i=1}^n x_i\leq\log k^* &\iff\notag\\
\sum_{i=1}^n x_i\leq\frac{1}{\theta_1-\theta_0}\log\left[\left(\frac{\theta_1}{\theta_0}\right)^n k^*\right]\equiv k'. & \quad\text{(constante)}
\end{align}
A desigualdade $\lambda(\v{x})\leq k^*$ foi simplificada e expressa como a desigualdade equivalente $\sum_{i=1}^n x_i\leq k'$.

A condição (i) do Lema de Neyman-Pearson requer
\begin{equation}
\alpha=\P_{\theta_0}\left(\sum_{i=1}^n X_i\leq k'\right).
\end{equation}
Observe que
\begin{equation}
\sum_{i=1}^n X_i\sim\text{gama}(n,\theta).
\end{equation}
Portanto, o teste mais poderoso de tamanho $\alpha$ de $\mathcal{H}_0:\theta=\theta_0$ \emph{vs.} $\mathcal{H}_1:\theta=\theta_1$ será aquele que rejeita $\mathcal{H}_0$ se $\sum_{i=1}^n x_i\leq k'$.

$k'$ corresponde ao $\alpha$-ésimo quantil de uma distribuição gama com parâmetros $n$ e $\theta_0$.
\end{example}
% Correspondência: 9 - Testes de Hipóteses 1.pdf, p. 40/45

% ---------- Exemplo 8 (Teste mais poderoso – Bernoulli) ----------
\begin{example}
Sejam $X_1,X_2,\ldots,X_n$ uma A.A. oriunda de uma população $X\sim\berndist{\theta}$, em que $\theta\in\{\theta_0,\theta_1\}$, com $\theta_0$ e $\theta_1$ números fixados.

Suponhamos aqui que queremos testar $\mathcal{H}_0:\theta=\theta_0$ \emph{vs.} $\mathcal{H}_1:\theta=\theta_1$, e que $\theta_0<\theta_1$.

Observe que
\begin{equation}
L_k=L(\theta_k\mid\v{x})=\theta_k^{\sum_{i=1}^n x_i}(1-\theta_k)^{n-\sum_{i=1}^n x_i}, \quad k=0,1.
\end{equation}
Assim,
\begin{align}
\lambda(\v{x}) &= \frac{L_0}{L_1}=\frac{\theta_0^{\sum_{i=1}^n x_i}(1-\theta_0)^{n-\sum_{i=1}^n x_i}}{\theta_1^{\sum_{i=1}^n x_i}(1-\theta_1)^{n-\sum_{i=1}^n x_i}}\notag\\
&= \left[\frac{\theta_0(1-\theta_1)}{(1-\theta_0)\theta_1}\right]^{\sum_{i=1}^n x_i}\left(\frac{1-\theta_0}{1-\theta_1}\right)^n.
\end{align}
Pelo Lema de Neyman-Pearson, o teste mais poderoso rejeita $\mathcal{H}_0$ se, e somente se, $\lambda(\v{x})\leq k^*$.

Portanto,
\begin{align}
\left[\frac{\theta_0(1-\theta_1)}{(1-\theta_0)\theta_1}\right]^{\sum_{i=1}^n x_i}\left(\frac{1-\theta_0}{1-\theta_1}\right)^n\leq k^* &\iff\notag\\
\sum_{i=1}^n x_i\log\left[\frac{\theta_0(1-\theta_1)}{(1-\theta_0)\theta_1}\right]+n\log\left(\frac{1-\theta_0}{1-\theta_1}\right)\leq\log k^* &\iff\notag\\
\sum_{i=1}^n x_i\geq\frac{\log\left[k^*\left(\dfrac{1-\theta_1}{1-\theta_0}\right)^n\right]}{\log\left[\dfrac{\theta_0(1-\theta_1)}{(1-\theta_0)\theta_1}\right]}\equiv k'. & \quad\text{(constante)}
\end{align}
Portanto, um teste mais poderoso é aquele que rejeita $\mathcal{H}_0$ se $\sum_{i=1}^n x_i$ for grande.

Para fins de ilustração, suponhamos que $\theta_0=1/4$, $\theta_1=3/4$ e $n=10$.

Devemos determinar $k'$ tal que
\begin{equation}
\alpha=\P_{\theta_0=1/4}(\text{rejeitar }\mathcal{H}_0)=\P_{\theta_0=1/4}\left(\sum_{i=1}^{10}X_i\geq k'\right)=\sum_{i=k'}^{10}\binom{10}{i}(1/4)^i(3/4)^{10-i}.
\end{equation}
Se $\alpha=0{,}0197$, então $k'=6$; se $\alpha=0{,}0781$, então $k'=5$.

Para $\alpha=0{,}05$, não há região crítica $\mathcal{C}^*$ e constante $k^*$ da forma como apresentado no Lema de Neyman-Pearson.

Para v.a.'s discretas, não é possível determinar $k^*$ e $\mathcal{C}^*$ que satisfaçam as condições (i) e (ii) para todo $0<\alpha<1$.

Na prática, é comum simplesmente mudar o tamanho do teste para algum $\alpha$ que possa garantir o referido lema. Contudo, é importante destacar que existe um teste (aleatorizado) mais poderoso.

Para o exemplo em questão, se $\alpha=0{,}05$, podemos considerar uma função crítica da forma
\begin{equation}
\Psi(\v{x})=
\begin{cases}
1, & \text{se }\sum_{i=1}^{10}x_i\geq6;\\
\gamma, & \text{se }\sum_{i=1}^{10}x_i=5;\\
0, & \text{se }\sum_{i=1}^{10}x_i\leq4.
\end{cases}
\end{equation}
Nesse caso,
\begin{equation}
\alpha=\E_{\theta_0=1/4}\left[\Psi(\v{X})\right]=1\cdot\P_{\theta_0=1/4}\left(\sum_{i=1}^n X_i\geq6\right)+\gamma\cdot\P_{\theta_0=1/4}\left(\sum_{i=1}^n X_i=5\right).
\end{equation}
Então,
\begin{equation}
0{,}05=0{,}0197+\gamma\cdot0{,}0584,
\end{equation}
o que implica que $\gamma=\dfrac{0{,}05-0{,}0197}{0{,}0584}=0{,}0303/0{,}0584$.

Logo,
\begin{equation}
\Psi(\v{x})=
\begin{cases}
1, & \text{se }\sum_{i=1}^{10}x_i\geq6,\\
0{,}0303/0{,}0584, & \text{se }\sum_{i=1}^{10}x_i=5\\
0, & \text{se }\sum_{i=1}^{10}x_i\leq4,
\end{cases}
\end{equation}
é o teste mais poderoso de tamanho $\alpha=0{,}05$.
\end{example}
% Correspondência: 9 - Testes de Hipóteses 1.pdf, p. 41/45

% ---------- Definição 10 (Nível descritivo / valor-p) ----------
\begin{definition}
Seja $\hat{p}$ o menor nível de significância para o qual a hipótese nula $\mathcal{H}_0$ é rejeitada, com base em uma determinada amostra observada. Essa quantidade é chamada de \uline{nível descritivo} (ou valor-$p$). Formalmente,
\begin{equation}
\hat{p}=\inf\{\alpha\in(0,1)\mid\v{x}\in\mathcal{C}_\alpha\},
\end{equation}
em que $\mathcal{C}_\alpha$ é uma região crítica construída com base em um teste de tamanho $\alpha$, e $\v{x}$ representa o vetor da amostra observada.
\end{definition}
% Correspondência: 9 - Testes de Hipóteses 1.pdf, p. 44/45

% ---------- Observação 6 (Interpretação do nível descritivo) ----------
\begin{note}
Quanto menor $\hat{p}$, mais evidência temos contra $\mathcal{H}_0$.
\end{note}
% Correspondência: 9 - Testes de Hipóteses 1.pdf, p. 44/45

% ---------- Exemplo 9 (Cont. – Ex. 7, nível descritivo) ----------
\begin{example}[Cont. – Exemplo 7]
Sejam $X_1,X_2,\ldots,X_n$ uma A.A. oriunda de uma população $X\sim\text{exp}(\theta)$, em que $\theta\in\{\theta_0,\theta_1\}$, com $\theta_0<\theta_1$ números fixados.

Estamos interessados em testar $\mathcal{H}_0:\theta=\theta_0$ \emph{vs.} $\mathcal{H}_1:\theta=\theta_1$.

Suponha que $n=20$, $\theta_0=1$ e $\theta_1=2$. Da amostra observada, foi verificado que $\sum_{i=1}^{20}x_i=6{,}18$.

Nesse caso, como o teste rejeita $\mathcal{H}_0$ para $\sum_{i=1}^{20}x_i\leq k'$, e $\sum_{i=1}^{20}X_i\sim\text{gama}(20,\theta)$, o nível descritivo será
\begin{equation}
\hat{p}=\P_{\theta=1}\left(\sum_{i=1}^{20}X_i\leq6{,}18\right)\approx0{,}0277.
\end{equation}
Portanto, ao nível de significância de 5\%, rejeitamos a hipótese nula de que $\theta=1$.
\end{example}
% Correspondência: 9 - Testes de Hipóteses 1.pdf, p. 45/45


\section{Testes de hipóteses II}
% !TeX root = main.tex
% !TeX spellcheck = pt_BR
% ===========================================================
% Testes de Hipóteses II: RVG e testes UMP
% Fonte: 10 - Testes de Hipóteses 2.pdf
% ===========================================================

% ---------- Introdução ----------
Até o momento, estudamos testes envolvendo hipóteses simples, o que não é usual em situações reais. Agora vamos dar ênfase a testes que envolvem hipóteses compostas.

Vamos trabalhar sob a suposição de que temos uma amostra aleatória oriunda de uma população com f.d.p. (ou f.p., embora bem menos)
\begin{equation}
f(\v{x}\mid\theta), \quad \theta\in\Theta.
\end{equation}
Nosso interesse será testar hipóteses $\mathcal{H}_0:\theta\in\Theta_0$ \emph{vs.} $\mathcal{H}_1:\theta\in\Theta_1$, em que $\Theta_0\subset\Theta$, $\Theta_1\subset\Theta$, $\Theta_0\cup\Theta_1=\Theta$ e $\Theta_0\cap\Theta_1=\emptyset$.
% Correspondência: 10 - Testes de Hipóteses 2.pdf, p. 3/95

% ---------- Definição 1 (RVG) ----------
\begin{definition}[Razão de verossimilhanças generalizada -- RVG]
Seja $L(\theta\mid\v{x})$ a função de verossimilhança para uma amostra $X_1,X_2,\ldots,X_n$, que tem f.d.p. (ou f.p.) conjunta $f(\v{x}\mid\theta)$, $\theta\in\Theta$. Chamamos de \uline{razão de verossimilhanças generalizada} a razão
\begin{equation}
\lambda=\lambda(\v{x})=\frac{\sup_{\theta\in\Theta_0}L(\theta\mid\v{x})}{\sup_{\theta\in\Theta}L(\theta\mid\v{x})}.
\end{equation}
\end{definition}
% Correspondência: 10 - Testes de Hipóteses 2.pdf, p. 5/95

% ---------- Texto: algumas ponderações sobre a RVG ----------
Observe que $\lambda$ é uma função da amostra observada $\v{x}$. Quando a RVG envolver as v.a.'s $\v{X}$, então a denotaremos por $\Lambda=\lambda(\v{X})$. Por vezes, podemos denotar $\Lambda_n$ ou $\lambda_n$, quando quisermos enfatizar o tamanho amostral.

Note também que $\Lambda$ é uma estatística, por ser uma função de $\v{X}$ que não depende de quantidades desconhecidas.

A partir da Definição 1, temos que $0\leq\lambda\leq1$. Uma vez que numerador e denominador não podem ser negativos, $\lambda\geq0$. Por outro lado, como $\Theta_0\subset\Theta$, então $\sup_{\theta\in\Theta_0}L(\theta\mid\v{x})\leq\sup_{\theta\in\Theta}L(\theta\mid\v{x})$, que implica $\lambda\leq1$.

Valores pequenos de $\lambda$ indicam evidências contra $\mathcal{H}_0$.
% Correspondência: 10 - Testes de Hipóteses 2.pdf, p. 6/95

Uma generalização natural do teste da razão de verossimilhanças, visto para testar hipóteses simples, seria
\begin{equation}
\lambda^*(\v{x})=\frac{\sup_{\theta\in\Theta_0}L(\theta\mid\v{x})}{\sup_{\theta\in\Theta_1}L(\theta\mid\v{x})}.
\end{equation}
Por ser mais simples de se calcular, o teste da razão de verossimilhanças generalizada é adotado na forma de $\lambda(\v{x})$, exibida na Definição 1.

Vale destacar que os testes $\lambda(\v{x})$ e $\lambda^*(\v{x})$ estão relacionados da forma
\begin{equation}
\lambda(\v{x})=\min\{1,\lambda^*(\v{x})\}.
\end{equation}
% Correspondência: 10 - Testes de Hipóteses 2.pdf, p. 7/95

% ---------- Observação 1 ----------
\begin{note}
	\leavevmode
	\begin{outline}[enumerate]
		\1 O parâmetro $\theta$ pode ser um escalar ou um vetor. Além disso, numerador e denominador de $\lambda(\v{x})$ correspondem aos EMV's de $\theta$ em $\Theta_0$ e $\Theta$, respectivamente, aplicados à função de verossimilhança.
	\end{outline}
\end{note}
% Correspondência: 10 - Testes de Hipóteses 2.pdf, p. 8/95

% ---------- Texto: princípio do TRVG e pontos negativos ----------
Valores $\lambda$ da estatística $\Lambda$ são usados para a tomada de decisão em um teste do tipo $\mathcal{H}_0:\theta\in\Theta_0$ \emph{vs.} $\mathcal{H}_1:\theta\in\Theta_1$. Para isso, utilizamos o princípio do teste da razão de verossimilhanças generalizada: rejeitar $\mathcal{H}_0$ se, e somente se, $\lambda\leq\lambda_0$, em que $\lambda_0$ é uma constante tal que $0<\lambda_0<1$. Essa constante é, geralmente, especificada fixando o tamanho do teste.

Alguns pontos negativos do TRVG:
\begin{outline}
\1 Às vezes é difícil determinar $\sup L(\theta\mid\v{x})$;
\1 Às vezes é difícil determinar a distribuição de $\Lambda$, fundamental para calcular o poder do teste.
\end{outline}
% Correspondência: 10 - Testes de Hipóteses 2.pdf, p. 9/95

% ---------- Observação 2 ----------
\begin{note}
	\leavevmode
	\begin{outline}[enumerate]
		\1 O desenvolvimento de testes da RVG geralmente começa de uma forma um tanto complicada, até chegar em resultados mais simples.
	\end{outline}
\end{note}
% Correspondência: 10 - Testes de Hipóteses 2.pdf, p. 11/95

% ---------- Observação 3 ----------
\begin{note}
	\leavevmode
	\begin{outline}[enumerate]
		\1 Por construção, testes da RVG dependem da amostra apenas através de estatísticas suficientes (mínimas).
	\end{outline}
\end{note}
% Correspondência: 10 - Testes de Hipóteses 2.pdf, p. 11/95

% ---------- Definição 2 (TUMP) ----------
\begin{definition}[Testes uniformemente mais poderosos -- TUMP]
Um teste $\Upsilon^*$ de $\mathcal{H}_0:\theta\in\Theta_0$ \emph{vs.} $\mathcal{H}_1:\theta\in\Theta_1$ é chamado de \uline{teste uniformemente mais poderoso} de tamanho $\alpha$ se, e somente se,
\begin{outline}
\1 (i) $\sup_{\theta\in\Theta_0}\pi_{\Upsilon^*}(\theta)=\alpha$;
\1 (ii) $\pi_{\Upsilon^*}(\theta)\geq\pi_\Upsilon(\theta)$, para todo $\theta\in\Theta_1$ e para qualquer teste $\Upsilon$ com tamanho menor do que ou igual a $\alpha$.
\end{outline}
\end{definition}
% Correspondência: 10 - Testes de Hipóteses 2.pdf, p. 13/95

% ---------- Observação 4 ----------
\begin{note}
	\leavevmode
	\begin{outline}[enumerate]
		\1 O advérbio \emph{uniformemente} está associado aos parâmetros $\theta\in\Theta_1$.
	\end{outline}
\end{note}
% Correspondência: 10 - Testes de Hipóteses 2.pdf, p. 13/95

% ---------- Texto de transição ----------
O Lema de Neyman-Pearson é fundamental para mostrarmos que um teste é UMP em casos como esse.

Há formas de estender o resultado apresentado para casos mais gerais. Os Teoremas 1 e 2 apresentados a seguir são fundamentais para isso.
% Correspondência: 10 - Testes de Hipóteses 2.pdf, p. 15/95

% ---------- Teorema 1 ----------
\begin{theorem}
Sejam $X_1,X_2,\ldots,X_n$ uma A.A. oriunda de uma população com f.d.p. $f(x\mid\theta)$, $\theta\in\Theta$, em que $\Theta$ é algum intervalo. Suponha que $f(x\mid\theta)$ pertence à família exponencial, ou seja,
\begin{equation}
f(x\mid\theta)=\exp\{\eta(\theta)T(x)-B(\theta)\}h(x).
\end{equation}
Seja $T^*(\v{x})\equiv\sum_{i=1}^n T(x_i)$.
\begin{outline}
\1 (i) Se $\eta(\theta)$ for uma função monótona crescente de $\theta$, e se existir uma constante $k^*$ tal que $\P_{\theta_0}(T^*(\v{X})>k^*)=\alpha$, então o teste $\Upsilon^*$ com região crítica $\mathcal{C}^*=\{\v{x}\in\mathbb{R}^n\mid T^*(\v{x})>k^*\}$ será um teste UMP de tamanho $\alpha$ para $\mathcal{H}_0:\theta\leq\theta_0$ \emph{vs.} $\mathcal{H}_1:\theta>\theta_0$ ou $\mathcal{H}_0:\theta=\theta_0$ \emph{vs.} $\mathcal{H}_1:\theta>\theta_0$;
\1 (ii) Se $\eta(\theta)$ for uma função monótona decrescente de $\theta$, e se existir uma constante $k^*$ tal que $\P_{\theta_0}(T^*(\v{X})<k^*)=\alpha$, então o teste $\Upsilon^*$ com região crítica $\mathcal{C}^*=\{\v{x}\in\mathbb{R}^n\mid T^*(\v{x})<k^*\}$ será um teste UMP de tamanho $\alpha$ para $\mathcal{H}_0:\theta\leq\theta_0$ \emph{vs.} $\mathcal{H}_1:\theta>\theta_0$ ou $\mathcal{H}_0:\theta=\theta_0$ \emph{vs.} $\mathcal{H}_1:\theta>\theta_0$.
\end{outline}
\end{theorem}
% Correspondência: 10 - Testes de Hipóteses 2.pdf, p. 16/95

% ---------- Definição 3 (Razão de verossimilhanças monótona) ----------
\begin{definition}[Razão de verossimilhanças monótona]
Uma família de f.d.p.'s $\{f(x\mid\theta)\mid\theta\in\Theta,\ \Theta\text{ intervalo}\}$ é dita ter \uline{razão de verossimilhanças monótona} se, de uma amostra $X_1,X_2,\ldots,X_n$, houver uma estatística $T\equiv T(\v{X})$ tal que $\dfrac{L(\theta'\mid\v{x})}{L(\theta''\mid\v{x})}$ seja uma função não crescente de $T(\v{x})$, para todo $\theta'<\theta''$, ou uma função não decrescente de $T(\v{x})$, para todo $\theta'<\theta''$.
\end{definition}
% Correspondência: 10 - Testes de Hipóteses 2.pdf, p. 18/95

% ---------- Observação 5 ----------
\begin{note}
	\leavevmode
	\begin{outline}[enumerate]
		\1 A razão de verossimilhanças monótona é apenas uma razão de duas funções de verossimilhança, e não uma razão de verossimilhanças generalizada.
	\end{outline}
\end{note}
% Correspondência: 10 - Testes de Hipóteses 2.pdf, p. 18/95

% ---------- Teorema 2 ----------
\begin{theorem}
Sejam $X_1,X_2,\ldots,X_n$ uma A.A. oriunda de uma população com f.d.p. $f(x\mid\theta)$, em que $\Theta$ é um intervalo. Suponha que a família de f.d.p.'s $\{f(x\mid\theta)\mid\theta\in\Theta\}$ tenha razão de verossimilhanças monótona na estatística $T(\v{X})$. Nesse caso:
\begin{outline}
\1 (i) Se a razão de verossimilhanças monótona for não decrescente em $T(\v{X})$, e se $k^*$ for tal que $\P_{\theta_0}(T(\v{X})<k^*)=\alpha$, então o teste correspondendo à região crítica $\mathcal{C}^*=\{\v{x}\in\mathbb{R}^n\mid T(\v{x})<k^*\}$ será UMP de tamanho $\alpha$ para $\mathcal{H}_0:\theta\leq\theta_0$ \emph{vs.} $\mathcal{H}_1:\theta>\theta_0$;
\1 (ii) Se a razão de verossimilhanças monótona for não crescente em $T(\v{x})$, e se $k^*$ for tal que $\P_{\theta_0}(T(\v{X})>k^*)=\alpha$, então o teste correspondendo à região crítica $\mathcal{C}^*=\{\v{x}\in\mathbb{R}^n\mid T(\v{x})>k^*\}$ será UMP de tamanho $\alpha$ para $\mathcal{H}_0:\theta\leq\theta_0$ \emph{vs.} $\mathcal{H}_1:\theta>\theta_0$.
\end{outline}
\end{theorem}
% Correspondência: 10 - Testes de Hipóteses 2.pdf, p. 21/95

% ---------- Observação 6 ----------
\begin{note}
	\leavevmode
	\begin{outline}[enumerate]
		\1 As hipóteses utilizadas nos Teoremas 1 e 2 são do tipo $\mathcal{H}_0:\theta\leq\theta_0$ \emph{vs.} $\mathcal{H}_1:\theta>\theta_0$. Contudo, esses mesmos teoremas podem ser reformulados para hipóteses na forma $\mathcal{H}_0:\theta\geq\theta_0$ \emph{vs.} $\mathcal{H}_1:\theta<\theta_0$, desde que os sinais de desigualdade das regiões críticas sejam invertidos.
	\end{outline}
\end{note}
% Correspondência: 10 - Testes de Hipóteses 2.pdf, p. 23/95

% ---------- Observação 7 ----------
\begin{note}
	\leavevmode
	\begin{outline}[enumerate]
		\1 Os referidos teoremas consideram apenas hipóteses unilaterais.
	\end{outline}
\end{note}
% Correspondência: 10 - Testes de Hipóteses 2.pdf, p. 23/95

% ---------- Texto: resumo da ópera ----------
Vimos que testes UMP existem para hipóteses unilaterais, quando a razão de verossimilhanças é monótona em alguma estatística.

Há vários problemas de testes de hipóteses em que não há um teste UMP.

Podemos pensar, portanto, em restringir a classe de testes, com o intuito de determinar um teste ótimo. Nesse caso, uma alternativa interessante corresponde aos testes não viesados.
% Correspondência: 10 - Testes de Hipóteses 2.pdf, p. 23/95

% ---------- Definição 4 (Teste não viesado) ----------
\begin{definition}
Um teste $\Upsilon$ da hipótese nula $\mathcal{H}_0:\theta\in\Theta_0$ \emph{vs.} a hipótese alternativa $\mathcal{H}_1:\theta\in\Theta_1$ é chamado de \uline{teste não viesado} se, e somente se,
\begin{equation}
\sup_{\theta\in\Theta_0}\pi_\Upsilon(\theta)\leq\inf_{\theta\in\Theta_1}\pi_\Upsilon(\theta).
\end{equation}
\end{definition}
% Correspondência: 10 - Testes de Hipóteses 2.pdf, p. 26/95

% ---------- Texto: testes não viesados ----------
Para testes não viesados, a probabilidade de rejeitar $\mathcal{H}_0$, quando essa hipótese for verdadeira, nunca será superior à probabilidade de rejeitá-la quando ela for falsa.

Se um teste não viesado for UMP, então ele será chamado de teste não viesado uniformemente mais poderoso.

Não entraremos em detalhe sobre essa classe de testes.
% Correspondência: 10 - Testes de Hipóteses 2.pdf, p. 26/95

% ---------- Contextualização: amostragem de uma população normal ----------
Suponhamos que a amostra aleatória $X_1,X_2,\ldots,X_n$ seja oriunda de uma população $X\sim\normdist{\mu,\sigma^2}$.

Vamos analisar testes de hipóteses para a média $\mu$ e para a variância $\sigma^2$.
% Correspondência: 10 - Testes de Hipóteses 2.pdf, p. 28/95

% ---------- Testes para a média: introdução ----------
O mais comum é estarmos interessados em testes unilaterais ou bilaterais.

No caso dos testes unilaterais, vamos abordar $\mathcal{H}_0:\mu\leq\mu_0$ \emph{vs.} $\mathcal{H}_1:\mu>\mu_0$.

Há dois casos a serem considerados:
\begin{outline}
\1 $\sigma^2$ conhecido;
\1 $\sigma^2$ desconhecido.
\end{outline}
% Correspondência: 10 - Testes de Hipóteses 2.pdf, p. 29/95

% ---------- Teorema/desenvolvimento: H0: mu <= mu0 vs H1: mu > mu0, sigma^2 conhecido ----------
\begin{theorem}[$\mathcal{H}_0:\mu\leq\mu_0$ \emph{vs.} $\mathcal{H}_1:\mu>\mu_0$ -- $\sigma^2$ conhecido]
Sejam $X_1,X_2,\ldots,X_n$ uma A.A. de uma população $X\sim\normdist{\mu,\sigma^2}$, com $\sigma^2$ conhecido. Nesse caso, o espaço paramétrico corresponderá ao conjunto dos reais. Além disso,
\begin{align}
f(x\mid\mu) &= \frac{1}{\sqrt{2\pi\sigma^2}}\exp\left\{-\frac{1}{2\sigma^2}(x-\mu)^2\right\}\notag\\
&= \frac{1}{\sqrt{2\pi\sigma^2}}\exp\left\{-\frac{1}{2\sigma^2}x^2+\frac{\mu}{\sigma^2}x-\frac{\mu^2}{2\sigma^2}\right\}\notag\\
&= \exp\left\{\frac{\mu}{\sigma^2}x-\frac{\mu^2}{2\sigma^2}\right\}\frac{\exp\left\{-\dfrac{x^2}{2\sigma^2}\right\}}{\sqrt{2\pi\sigma^2}}\notag\\
&= \exp\{\eta(\mu)T(x)-B(\mu)\}h(x).
\end{align}
Como $\eta(\mu)=\mu/\sigma^2$ é uma função crescente de $\mu$, então, pelo Teorema 1, o teste UMP é aquele que rejeita $\mathcal{H}_0$ se, e somente se, $\sum_{i=1}^n x_i>k^*$, em que a constante $k^*$ é dada pela solução de
\begin{equation}
\alpha=\P_{\mu_0}\left(\sum_{i=1}^n X_i>k^*\right)=1-\Phi\left(\frac{k^*-n\mu_0}{\sqrt{n}\sigma}\right).
\end{equation}
Portanto, $\dfrac{k^*-n\mu_0}{\sqrt{n}\sigma}=z_{1-\alpha}$, em que $z_{1-\alpha}$ corresponde ao $(1-\alpha)$-ésimo quantil da distribuição normal padrão.

Desse modo, o teste UMP rejeita $\mathcal{H}_0$ se $\sum_{i=1}^n x_i>n\mu_0+\sqrt{n}\sigma z_{1-\alpha}$ ou, de forma equivalente, rejeita $\mathcal{H}_0$ se $\bar{x}>\mu_0+\dfrac{\sigma}{\sqrt{n}}z_{1-\alpha}$.
\end{theorem}
% Correspondência: 10 - Testes de Hipóteses 2.pdf, p. 31-32/95

% ---------- Teorema/desenvolvimento: H0: mu <= mu0 vs H1: mu > mu0, sigma^2 desconhecido ----------
\begin{theorem}[$\mathcal{H}_0:\mu\leq\mu_0$ \emph{vs.} $\mathcal{H}_1:\mu>\mu_0$ -- $\sigma^2$ desconhecido]
Sejam $X_1,X_2,\ldots,X_n$ uma A.A. de uma população $X\sim\normdist{\mu,\sigma^2}$, com $\sigma^2$ desconhecido. Nesse caso, $\Theta=\{(\mu,\sigma^2)\in\mathbb{R}^2\mid-\infty<\mu<\infty;\sigma^2>0\}$ e $\Theta_0=\{(\mu,\sigma^2)\in\mathbb{R}^2\mid\mu\leq\mu_0;\sigma^2>0\}$.

No caso em que $(\mu,\sigma^2)\in\Theta$, temos que os EMV's de $\mu$ e $\sigma^2$ correspondem a $\hat{\mu}=\bar{X}$ e $\hat{\sigma}^2=\sum_{i=1}^n(X_i-\bar{X})^2/n$, respectivamente. Assim,
\begin{align}
\sup_{(\mu,\sigma^2)\in\Theta}L(\mu,\sigma^2\mid\v{x}) &= L(\hat{\mu},\hat{\sigma}^2\mid\v{x}) = \frac{1}{(2\pi\hat{\sigma}^2)^{n/2}}\exp\left[-\frac{1}{2\hat{\sigma}^2}\sum_{i=1}^n(x_i-\hat{\mu})^2\right]\notag\\
&= \left[\frac{n}{2\pi\sum_{i=1}^n(x_i-\bar{x})^2}\right]^{n/2}\exp\left(-\frac{n}{2}\right).
\end{align}
Quando $(\mu,\sigma^2)\in\Theta_0$, pode-se mostrar (mostre!) que os EMV's dos parâmetros serão
\begin{equation}
\tilde{\mu} = \begin{cases}
\bar{X}, & \text{se } \bar{X}\leq\mu_0;\\
\mu_0, & \text{se } \bar{X}>\mu_0
\end{cases}
\quad \text{e} \quad
\tilde{\sigma}^2 = \begin{cases}
\dfrac{1}{n}\sum_{i=1}^n\left(X_i-\bar{X}\right)^2, & \text{se } \bar{X}\leq\mu_0;\\
\dfrac{1}{n}\sum_{i=1}^n(X_i-\mu_0)^2, & \text{se } \bar{X}>\mu_0.
\end{cases}
\end{equation}
Desse modo,
\begin{align}
\sup_{(\mu,\sigma^2)\in\Theta_0}L(\mu,\sigma^2\mid\v{x}) &= L(\tilde{\mu},\tilde{\sigma}^2\mid\v{x}) = \frac{1}{(2\pi\tilde{\sigma}^2)^{n/2}}\exp\left[-\frac{1}{2\tilde{\sigma}^2}\sum_{i=1}^n(x_i-\tilde{\mu})^2\right]\notag\\
&= \begin{cases}
\left[\dfrac{n}{2\pi\sum_{i=1}^n(x_i-\bar{x})^2}\right]^{n/2}\exp\left(-\dfrac{n}{2}\right), & \text{se } \bar{x}\leq\mu_0;\\[2ex]
\left[\dfrac{n}{2\pi\sum_{i=1}^n(x_i-\mu_0)^2}\right]^{n/2}\exp\left(-\dfrac{n}{2}\right), & \text{se } \bar{x}>\mu_0.
\end{cases}
\end{align}
A razão de verossimilhanças será
\begin{align}
\lambda(\v{x}) &= \frac{\sup_{(\mu,\sigma^2)\in\Theta_0}L(\mu,\sigma^2\mid\v{x})}{\sup_{(\mu,\sigma^2)\in\Theta}L(\mu,\sigma^2\mid\v{x})}\notag\\
&= \begin{cases}
1, & \text{se } \bar{x}\leq\mu_0;\\[1ex]
\left[\dfrac{\sum_{i=1}^n(x_i-\bar{x})^2}{\sum_{i=1}^n(x_i-\mu_0)^2}\right]^{n/2}, & \text{se } \bar{x}>\mu_0.
\end{cases}
\end{align}
Antes de continuarmos, observe que
\begin{align}
\frac{\sum_{i=1}^n(x_i-\bar{x})^2}{\sum_{i=1}^n(x_i-\mu_0)^2} &= \frac{\sum_{i=1}^n(x_i-\bar{x})^2}{\sum_{i=1}^n(x_i-\bar{x}+\bar{x}-\mu_0)^2} = \frac{\sum_{i=1}^n(x_i-\bar{x})^2}{\sum_{i=1}^n(x_i-\bar{x})^2+n(\bar{x}-\mu_0)^2}\notag\\
&= \frac{1}{1+n\dfrac{(\bar{x}-\mu_0)^2}{\sum_{i=1}^n(x_i-\bar{x})^2}} = \frac{1}{1+n\dfrac{(\bar{x}-\mu_0)^2}{(n-1)s^2}} = \frac{1}{1+\dfrac{1}{n-1}\dfrac{(\bar{x}-\mu_0)^2}{s^2/n}}\notag\\
&= \left[1+\frac{1}{n-1}\frac{(\bar{x}-\mu_0)^2}{s^2/n}\right]^{-1} = \left(1+\frac{t^2}{n-1}\right)^{-1},
\end{align}
em que $t=\dfrac{\bar{x}-\mu_0}{s/\sqrt{n}}$. Logo,
\begin{equation}
\lambda(\v{x}) = \begin{cases}
1, & \text{se } \bar{x}\leq\mu_0;\\[1ex]
\left(1+\dfrac{t^2}{n-1}\right)^{-n/2}, & \text{se } \bar{x}>\mu_0.
\end{cases}
\end{equation}
Desse modo, para $0<\lambda_0<1$, rejeitamos $\mathcal{H}_0$ se, e somente se,
\begin{equation}
\lambda(\v{x})\leq\lambda_0 \iff t\geq k.
\end{equation}
Denotemos $T=\dfrac{\bar{X}-\mu_0}{S/\sqrt{n}}$. Sabemos que, se $\mu=\mu_0$, então $T\sim\tdist{n-1}$. Portanto, o valor de $k$ é tal que
\begin{equation}
\alpha=\P_{\mu_0}(T\geq k)=\P_{\mu_0}\left(T\geq t_{n-1,1-\alpha}\right),
\end{equation}
em que $t_{n-1,1-\alpha}$ corresponde ao $(1-\alpha)$-ésimo quantil de uma $\tdist{n-1}$.

O teste rejeita $\mathcal{H}_0$ se, e somente se,
\begin{equation}
\frac{\bar{x}-\mu_0}{s/\sqrt{n}}\geq t_{n-1,1-\alpha},
\end{equation}
ou, de forma equivalente,
\begin{equation}
\bar{x}\geq\mu_0+\frac{s}{\sqrt{n}}t_{n-1,1-\alpha}.
\end{equation}
\end{theorem}
% Correspondência: 10 - Testes de Hipóteses 2.pdf, p. 33-39/95

% ---------- Teorema/desenvolvimento: H0: mu = mu0 vs H1: mu != mu0, sigma^2 desconhecido ----------
\begin{theorem}[$\mathcal{H}_0:\mu=\mu_0$ \emph{vs.} $\mathcal{H}_1:\mu\neq\mu_0$ -- $\sigma^2$ desconhecido]
Sejam $X_1,X_2,\ldots,X_n$ uma A.A. de uma população $X\sim\normdist{\mu,\sigma^2}$, com $\sigma^2$ desconhecido. No caso em que $\sigma^2$ for desconhecido, o espaço paramétrico irrestrito e o espaço paramétrico sob $\mathcal{H}_0$ serão $\Theta=\{(\mu,\sigma^2)\in\mathbb{R}^2\mid-\infty<\mu<\infty;\sigma^2>0\}$ e $\Theta_0=\{(\mu,\sigma^2)\in\mathbb{R}^2\mid\mu=\mu_0;\sigma^2>0\}$.

Já vimos que
\begin{equation}
\sup_{(\mu,\sigma^2)\in\Theta}L(\mu,\sigma^2\mid\v{x}) = \left[\frac{n}{2\pi\sum_{i=1}^n(x_i-\bar{x})^2}\right]^{n/2}\exp\left(-\frac{n}{2}\right).
\end{equation}
Quando $(\mu,\sigma^2)\in\Theta_0$, não é difícil verificar (verifique!) que os EMV's dos parâmetros são, respectivamente,
\begin{equation}
\tilde{\mu}=\mu_0 \quad \text{e} \quad \tilde{\sigma}^2=\frac{1}{n}\sum_{i=1}^n(X_i-\mu_0)^2.
\end{equation}
Logo,
\begin{align}
\sup_{(\mu,\sigma^2)\in\Theta_0}L(\mu,\sigma^2\mid\v{x}) &= L(\tilde{\mu},\tilde{\sigma}^2\mid\v{x}) = \frac{1}{(2\pi\tilde{\sigma}^2)^{n/2}}\exp\left[-\frac{1}{2\tilde{\sigma}^2}\sum_{i=1}^n(x_i-\tilde{\mu})^2\right]\notag\\
&= \left[\frac{n}{2\pi\sum_{i=1}^n(x_i-\mu_0)^2}\right]^{n/2}\exp\left(-\frac{n}{2}\right).
\end{align}
De forma análoga ao que vimos anteriormente,
\begin{align}
\lambda(\v{x}) &= \frac{\sup_{(\mu,\sigma^2)\in\Theta_0}L(\mu,\sigma^2\mid\v{x})}{\sup_{(\mu,\sigma^2)\in\Theta}L(\mu,\sigma^2\mid\v{x})} = \left[\frac{\sum_{i=1}^n(x_i-\bar{x})^2}{\sum_{i=1}^n(x_i-\mu_0)^2}\right]^{n/2}\notag\\
&= \left(1+\frac{t^2}{n-1}\right)^{-n/2},
\end{align}
em que $t=\dfrac{\bar{x}-\mu_0}{s/\sqrt{n}}$.

Assim, para $0<\lambda_0<1$,
\begin{align}
\left(1+\frac{t^2}{n-1}\right)^{-n/2}\leq\lambda_0 &\iff\notag\\
|t|\geq\sqrt{(n-1)\left(\lambda_0^{-2/n}-1\right)}\equiv k &\iff\notag\\
t\leq-k \quad \text{ou} \quad t\geq k.
\end{align}
Denotemos $T=\dfrac{\bar{X}-\mu_0}{S/\sqrt{n}}$. Sabemos que, sob $\mathcal{H}_0$, $T\sim\tdist{n-1}$.

O valor de $k$ será solução de
\begin{align}
\alpha &= \P_{\mu_0}\left([T\leq-k]\text{ ou }[T\geq k]\right) = \P_{\mu_0}(T\leq-k)+\P_{\mu_0}(T\geq k)\notag\\
&= 2\P_{\mu_0}(T\geq k).
\end{align}
Portanto, $k=t_{n-1,1-\alpha/2}$.

Desse modo, o teste da RVG rejeita $\mathcal{H}_0$ se, e somente se,
\begin{equation}
t\leq-t_{n-1,1-\alpha/2} \quad \text{ou} \quad t\geq t_{n-1,1-\alpha/2},
\end{equation}
ou seja,
\begin{equation}
\bar{x}\leq\mu_0-t_{n-1,1-\alpha/2}\frac{s}{\sqrt{n}} \quad \text{ou} \quad \bar{x}\geq\mu_0+t_{n-1,1-\alpha/2}\frac{s}{\sqrt{n}}.
\end{equation}
Vale destacar que esse teste não é UMP. Por outro lado, pode-se mostrar que é um teste não viesado uniformemente mais poderoso.
\end{theorem}
% Correspondência: 10 - Testes de Hipóteses 2.pdf, p. 41-47/95

% ---------- Testes para a variância: introdução ----------
Assim como feito para testes para a média, vamos abordar testes unilaterais do tipo $\mathcal{H}_0:\sigma^2\leq\sigma_0^2$ \emph{vs.} $\mathcal{H}_1:\sigma^2>\sigma_0^2$, bem como testes bilaterais.

Há dois casos a serem considerados:
\begin{outline}
\1 $\mu$ conhecido;
\1 $\mu$ desconhecido.
\end{outline}
% Correspondência: 10 - Testes de Hipóteses 2.pdf, p. 48/95

% ---------- Teorema/desenvolvimento: H0: sigma^2 <= sigma0^2 vs H1: sigma^2 > sigma0^2, mu conhecido ----------
\begin{theorem}[$\mathcal{H}_0:\sigma^2\leq\sigma_0^2$ \emph{vs.} $\mathcal{H}_1:\sigma^2>\sigma_0^2$ -- $\mu$ conhecido]
Sejam $X_1,X_2,\ldots,X_n$ uma A.A. de uma população $X\sim\normdist{\mu,\sigma^2}$, com $\mu$ conhecido. Nesse caso, o espaço paramétrico será um intervalo (conjunto dos reais positivos). Além disso,
\begin{align}
f(x\mid\mu) &= \frac{1}{\sqrt{2\pi\sigma^2}}\exp\left\{-\frac{1}{2\sigma^2}(x-\mu)^2\right\}\notag\\
&= \exp\left\{-\frac{1}{2\sigma^2}(x-\mu)^2-\frac{1}{2}\log\sigma^2\right\}\frac{1}{\sqrt{2\pi}}\notag\\
&= \exp\{\eta(\sigma^2)T(x)-B(\sigma^2)\}h(x).
\end{align}
Seja $R=\sum_{i=1}^n\left(\dfrac{X_i-\mu}{\sigma_0}\right)^2$. Já sabemos que, se a verdadeira variância for $\sigma_0^2$, então $R\sim\chidist{n}$.

Observe que $\eta(\sigma^2)=-1/(2\sigma^2)$ é uma função crescente de $\sigma^2$. Pelo Teorema 1, o teste UMP rejeita $\mathcal{H}_0$ se, e somente se, $\sum_{i=1}^n(x_i-\mu)^2>k^*$, em que $k^*$ é solução de
\begin{align}
\alpha=\P_{\sigma_0^2}\left(\sum_{i=1}^n(X_i-\mu)^2>k^*\right) &= \P_{\sigma_0^2}\left(R>k^*/\sigma_0^2\right) = \P_{\sigma_0^2}\left(R>\chi^2_{n,1-\alpha}\right).
\end{align}
Portanto, $\dfrac{k^*}{\sigma_0^2}=\chi^2_{n,1-\alpha}$.

Assim, o teste UMP rejeita $\mathcal{H}_0$ se, e somente se, $\sum_{i=1}^n(x_i-\mu)^2>\sigma_0^2\chi^2_{n,1-\alpha}$.
\end{theorem}
% Correspondência: 10 - Testes de Hipóteses 2.pdf, p. 49-50/95

% ---------- Teorema/desenvolvimento: H0: sigma^2 = sigma0^2 vs H1: sigma^2 != sigma0^2, mu desconhecido ----------
\begin{theorem}[$\mathcal{H}_0:\sigma^2=\sigma_0^2$ \emph{vs.} $\mathcal{H}_1:\sigma^2\neq\sigma_0^2$ -- $\mu$ desconhecido]
Sejam $X_1,X_2,\ldots,X_n$ uma A.A. de uma população $X\sim\normdist{\mu,\sigma^2}$, com $\mu$ desconhecido. Nesse caso, $\Theta=\{(\mu,\sigma^2)\in\mathbb{R}^2\mid-\infty<\mu<\infty;\sigma^2>0\}$ e $\Theta_0=\{(\mu,\sigma^2)\in\mathbb{R}^2\mid-\infty<\mu<\infty;\sigma^2=\sigma_0^2\}$.

Já vimos que
\begin{equation}
\sup_{(\mu,\sigma^2)\in\Theta}L(\mu,\sigma^2\mid\v{x}) = \left(\frac{1}{2\pi\hat{\sigma}^2}\right)^{n/2}\exp\left(-\frac{n}{2}\right).
\end{equation}
Quando $(\mu,\sigma^2)\in\Theta_0$, não é difícil verificar (verifique!) que os EMV's dos parâmetros são, respectivamente,
\begin{equation}
\tilde{\mu}=\bar{X} \quad \text{e} \quad \tilde{\sigma}^2=\sigma_0^2.
\end{equation}
Logo,
\begin{align}
\sup_{(\mu,\sigma^2)\in\Theta_0}L(\mu,\sigma^2\mid\v{x}) &= L(\tilde{\mu},\tilde{\sigma}^2\mid\v{x}) = \frac{1}{(2\pi\tilde{\sigma}^2)^{n/2}}\exp\left[-\frac{1}{2\tilde{\sigma}^2}\sum_{i=1}^n(x_i-\tilde{\mu})^2\right]\notag\\
&= \frac{1}{(2\pi\sigma_0^2)^{n/2}}\exp\left[-\frac{1}{2\sigma_0^2}\sum_{i=1}^n(x_i-\bar{x})^2\right] = \frac{1}{(2\pi\sigma_0^2)^{n/2}}\exp\left(-\frac{n\hat{\sigma}^2}{2\sigma_0^2}\right).
\end{align}
Assim,
\begin{align}
\lambda(\v{x}) &= \frac{\sup_{(\mu,\sigma^2)\in\Theta_0}L(\mu,\sigma^2\mid\v{x})}{\sup_{(\mu,\sigma^2)\in\Theta}L(\mu,\sigma^2\mid\v{x})} = \left(\frac{\hat{\sigma}^2}{\sigma_0^2}\right)^{n/2}\exp\left(-\frac{n}{2}\frac{\hat{\sigma}^2}{\sigma_0^2}+\frac{n}{2}\right)\notag\\
&\equiv y^{n/2}\exp\left[-\frac{n}{2}(y-1)\right],
\end{align}
em que $y=\dfrac{\hat{\sigma}^2}{\sigma_0^2}$.

Assim, para $0<\lambda_0<1$,
\begin{align}
y^{n/2}\exp\left[-\frac{n}{2}(y-1)\right]\leq\lambda_0 &\iff\notag\\
y\leq k_1 \quad \text{ou} \quad y\geq k_2, & 
\end{align}
em que $y=\dfrac{\hat{\sigma}^2}{\sigma_0^2}$.

Seja $V=\dfrac{\sum_{i=1}^n\left(X_i-\bar{X}\right)^2}{\sigma_0^2}$. Nesse caso, sob $\mathcal{H}_0$, $V\sim\chidist{n-1}$. Além disso, note que $\dfrac{\hat{\sigma}^2}{\sigma_0^2}=\dfrac{V}{n}$.

Assim, os valores de $k_1$ e $k_2$ são soluções de
\begin{align}
\alpha &= \P_{\sigma_0^2}\left(\left[\frac{\hat{\sigma}^2}{\sigma_0^2}\leq k_1\right]\text{ ou }\left[\frac{\hat{\sigma}^2}{\sigma_0^2}\geq k_2\right]\right) = \P_{\sigma_0^2}\left(\left[\frac{V}{n}\leq k_1\right]\text{ ou }\left[\frac{V}{n}\geq k_2\right]\right)\notag\\
&= \P_{\sigma_0^2}(V\leq nk_1)+\P_{\mu_0}(V\geq nk_2)\notag\\
&= \P_{\sigma_0^2}\left(V\leq\chi^2_{n-1,\alpha/2}\right)+\P_{\mu_0}\left(V\geq\chi^2_{n-1,1-\alpha/2}\right) \quad \text{(uma possibilidade).}
\end{align}
Nesse caso mais simples (que não é o melhor!), temos
\begin{equation}
k_1=\frac{\chi^2_{n-1,\alpha/2}}{n} \quad \text{e} \quad k_2=\frac{\chi^2_{n-1,1-\alpha/2}}{n}.
\end{equation}
Desse modo, o teste rejeita $\mathcal{H}_0$ se, e somente se,
\begin{equation}
\frac{\hat{\sigma}^2}{\sigma_0^2}\leq\frac{\chi^2_{n-1,\alpha/2}}{n} \quad \text{ou} \quad \frac{\hat{\sigma}^2}{\sigma_0^2}\geq\frac{\chi^2_{n-1,1-\alpha/2}}{n},
\end{equation}
em que $\chi^2_{r,\beta}$ corresponde ao $\beta$-ésimo quantil de uma distribuição $\chidist{r}$.
\end{theorem}
% Correspondência: 10 - Testes de Hipóteses 2.pdf, p. 53-59/95

% ---------- Texto: distribuição assintótica do TRVG ----------
À medida que o número de parâmetros do modelo aumenta, torna-se mais complicada a tarefa de estabelecer testes da RVG.

Por vezes, essa tarefa não é tão simples nem mesmo em situações uniparamétricas.

Uma forma de contornar esse problema é através de aproximação.
% Correspondência: 10 - Testes de Hipóteses 2.pdf, p. 61/95

% ---------- Teorema 3 ----------
\begin{theorem}
Sejam $X_1,X_2,\ldots,X_n$ uma A.A. de uma população com f.d.p. (ou f.p.) $f(x\mid\theta)$, em que $\theta$ pode ser escalar ou vetor. Sob certas condições de regularidade, se $\theta\in\Theta_0$, então
\begin{equation}
-2\log\Lambda_n \toD \chidist{r},
\end{equation}
quando $n\to\infty$. O número de graus de liberdade $r$ corresponde ao número de restrições em $\Theta_0$.
\end{theorem}
% Correspondência: 10 - Testes de Hipóteses 2.pdf, p. 62/95

% ---------- Observação 8 ----------
\begin{note}
	\leavevmode
	\begin{outline}[enumerate]
		\1 O número de graus de liberdade $r$ também pode ser interpretado como a dimensão de $\Theta$ menos a dimensão de $\Theta_0$.
	\end{outline}
\end{note}
% Correspondência: 10 - Testes de Hipóteses 2.pdf, p. 62/95

% ---------- Texto: princípio da RVG no caso assintótico ----------
Lembrando que $\Lambda_n$ é uma v.a. que assume valores do tipo
\begin{equation}
\lambda_n=\frac{\sup_{\theta\in\Theta_0}L(\theta\mid\v{x})}{\sup_{\theta\in\Theta}L(\theta\mid\v{x})},
\end{equation}
que corresponde à RVG para uma amostra observada $\v{x}=(x_1,x_2,\ldots,x_n)$.

O princípio da RVG sugere que rejeitemos $\mathcal{H}_0$ para valores pequenos de $\lambda_n$. Portanto, no caso assintótico, o teste deve rejeitar a hipótese nula para valores grandes de $-2\log\lambda_n$.
% Correspondência: 10 - Testes de Hipóteses 2.pdf, p. 63/95

Levando em consideração a distribuição assintótica quando $\mathcal{H}_0$ for verdadeira, um teste com tamanho aproximadamente $\alpha$ rejeita $\mathcal{H}_0$ se, e somente se,
\begin{equation}
-2\log\lambda_n\geq\chi^2_{r,1-\alpha},
\end{equation}
em que $\chi^2_{r,1-\alpha}$ corresponde ao $(1-\alpha)$-ésimo quantil de uma distribuição $\chidist{r}$.
% Correspondência: 10 - Testes de Hipóteses 2.pdf, p. 64/95


